
\documentclass{article}
\usepackage{colortbl}
\usepackage{makecell}
\usepackage{multirow}
\usepackage{supertabular}

\begin{document}

\newcounter{utterance}

\centering \large Interaction Transcript for game `ta\_mastermind', experiment `Mastermind{-}v0', episode 6 with qwen3.5{-}9b.
\vspace{24pt}

{ \footnotesize  \setcounter{utterance}{1}
\setlength{\tabcolsep}{0pt}
\begin{supertabular}{c@{$\;$}|p{.15\linewidth}@{}p{.15\linewidth}p{.15\linewidth}p{.15\linewidth}p{.15\linewidth}p{.15\linewidth}}
    \# & \multicolumn{2}{c}{Player} && \multicolumn{2}{c}{Game Master} \\
    \hline

    \theutterance \stepcounter{utterance}  
    & & & \multicolumn{4}{p{0.6\linewidth}}{
        \cellcolor[rgb]{0.9,0.9,0.9}{
            \makecell[{{p{\linewidth}}}]{
                \texttt{\tiny{[P1$\langle$GM]}}
                \texttt{You are a competitive game player. Make sure you read the game instructions carefully, and always follow the required format.} \\
\\ 
\texttt{{[}GAME{]} You are playing Mastermind.} \\
\texttt{You need to find the code that is 4 digits long, each digit from 1 to 6, with no duplicates.} \\
\texttt{In your response, you can submit your guess in the following format: '{[}2 1 4 5{]}'.} \\
\texttt{After each guess, you will receive feedback in the form of black and white pegs.} \\
\texttt{A black peg indicates a correct digit in the correct position, while a white peg indicates a correct digit in the wrong position.} \\
\texttt{You have 20 turns to guess the code.} \\
\\ 
            }
        }
    }
     \\ \\

    \theutterance \stepcounter{utterance}  
    & \multicolumn{4}{p{0.6\linewidth}}{
        \cellcolor[rgb]{0.561,0.737,0.561}{
            \makecell[{{p{\linewidth}}}]{
                \texttt{\tiny{[P1$\rangle$GM]}}
                \texttt{{[}1 2 3 4{]}} \\
            }
        }
    }
    & & \\ \\

    \theutterance \stepcounter{utterance}  
    & & & \multicolumn{4}{p{0.6\linewidth}}{
        \cellcolor[rgb]{0.9,0.9,0.9}{
            \makecell[{{p{\linewidth}}}]{
                \texttt{\tiny{[P1$\langle$GM]}}
                \texttt{{[}GAME{]} Submitted {[}1 2 3 4{]}. Feedback: 0 black peg(s), 2 white peg(s).} \\
            }
        }
    }
     \\ \\

    \theutterance \stepcounter{utterance}  
    & \multicolumn{4}{p{0.6\linewidth}}{
        \cellcolor[rgb]{0.561,0.737,0.561}{
            \makecell[{{p{\linewidth}}}]{
                \texttt{\tiny{[P1$\rangle$GM]}}
                \texttt{Based on the feedback (0 black, 2 white), I know that two of my digits \{1, 2, 3, 4\} are in the code but in the wrong positions, and the other two are not in the code at all.} \\
\\ 
\texttt{Since there are no duplicates allowed, I need to eliminate the incorrect digits and test new combinations. I'll try to isolate which digits are correct by swapping positions or introducing new numbers from the remaining set \{5, 6\}.} \\
\\ 
\texttt{Let's guess: {[}2 1 5 6{]}} \\
\texttt{{-} If 1 and 2 are the correct digits:} \\
\texttt{  {-} Previous guess: 1(pos1), 2(pos2). New guess: 2(pos1), 1(pos2). Both are swapped. If they were the correct digits, the white peg count should remain 2 (just swapped positions).} \\
\texttt{  {-} 5 and 6 are new. If either is correct, the white peg count would increase.} \\
\texttt{{-} This guess helps determine if 1 and 2 are the pair, or if one of them is wrong and 5 or 6 is right.} \\
\\ 
\texttt{{[}2 1 5 6{]}} \\
            }
        }
    }
    & & \\ \\

    \theutterance \stepcounter{utterance}  
    & & & \multicolumn{4}{p{0.6\linewidth}}{
        \cellcolor[rgb]{0.9,0.9,0.9}{
            \makecell[{{p{\linewidth}}}]{
                \texttt{\tiny{[P1$\langle$GM]}}
                \texttt{{[}GAME{]} Submitted {[}2 1 5 6{]}. Feedback: 1 black peg(s), 2 white peg(s).} \\
            }
        }
    }
     \\ \\

    \theutterance \stepcounter{utterance}  
    & \multicolumn{4}{p{0.6\linewidth}}{
        \cellcolor[rgb]{0.561,0.737,0.561}{
            \makecell[{{p{\linewidth}}}]{
                \texttt{\tiny{[P1$\rangle$GM]}}
                \texttt{Based on the feedback from my previous guess `{[}2 1 5 6{]}` (1 black, 2 white), let's analyze the situation:} \\
\\ 
\texttt{1.  **From Guess 1 `{[}1 2 3 4{]}` (0B, 2W):** Two digits from \{1, 2, 3, 4\} are in the code but in wrong positions. Two are not in the code.} \\
\texttt{2.  **From Guess 2 `{[}2 1 5 6{]}` (1B, 2W):**} \\
\texttt{    *   We have 1 Black peg. This means one digit is correct and in the correct position.} \\
\texttt{    *   We have 2 White pegs. This means two digits are correct but in the wrong position.} \\
\texttt{    *   Total correct digits identified so far: 1 (Black) + 2 (White) = 3 digits.} \\
\texttt{    *   Since the code has 4 digits, there is exactly **1 digit left to find**.} \\
\\ 
\texttt{Let's look at the overlap between the two guesses:} \\
\texttt{*   Guess 1: `1 2 3 4`} \\
\texttt{*   Guess 2: `2 1 5 6`} \\
\texttt{*   Common digits: `1` and `2`.} \\
\\ 
\texttt{In Guess 1, `1` was at pos 1 and `2` was at pos 2. Neither was a Black peg (0B).} \\
\texttt{In Guess 2, `2` is at pos 1 and `1` is at pos 2. We got 1 Black peg.} \\
\texttt{This implies that either `2` is now in the correct position (pos 1) or `1` is now in the correct position (pos 2).} \\
\\ 
\texttt{*   **Scenario A:** `2` is the Black peg (Code has `2` at pos 1).} \\
\texttt{    *   If `2` is correct at pos 1, then in Guess 1 (`1 2 3 4`), `2` was at pos 2 (wrong pos), which fits the 0B result.} \\
\texttt{    *   The 2 White pegs in Guess 2 must come from `1`, `5`, or `6`. Since `2` is the Black peg, the other 2 correct digits are among `\{1, 5, 6\}`.} \\
\texttt{    *   In Guess 1, we had 2 White pegs from `\{1, 2, 3, 4\}`. If `2` is one of them, the other must be `1`, `3`, or `4`.} \\
\texttt{    *   If `2` is the Black peg, `1` cannot be a White peg in Guess 2 if it's the *same* `1` from Guess 1? Wait, let's re{-}evaluate the set of correct digits.} \\
\texttt{    *   Let's assume the correct digits are `\{2, X, Y\}`.} \\
\texttt{    *   If `2` is the Black peg (pos 1), then `1` in Guess 2 is in the wrong position (pos 2). So `1` would be a White peg. That accounts for 1 White. We need 1 more White from `\{5, 6\}`.} \\
\texttt{    *   So the code would be `2 \_ \_ \_` with one of `\{1, 5, 6\}` being the second correct digit and one of `\{3, 4\}` being the third? No, we only have 3 correct digits total (1B + 2W).} \\
\texttt{    *   If `2` is the Black peg, the other two correct digits must be from the set that provided the White pegs.} \\
\texttt{    *   In Guess 1 (`1 2 3 4`), if `2` is correct (but at pos 2, so it's a White peg there), we need one more White peg from `\{1, 3, 4\}`.} \\
\texttt{    *   In Guess 2 (`2 1 5 6`), `2` is Black. We need 2 White pegs. One is `1` (at pos 2, wrong). The other must be `5` or `6`.} \\
\texttt{    *   So the correct digits are likely `\{2, 1, 5\}` or `\{2, 1, 6\}` or `\{2, 5, 6\}` (if 1 is wrong).} \\
\texttt{    *   If the set is `\{2, 1, 5\}`:} \\
\texttt{        *   Guess 1 (`1 2 3 4`): `1` (pos 1, wrong), `2` (pos 2, wrong). 2 Whites. Matches.} \\
\texttt{        *   Guess 2 (`2 1 5 6`): `2` (pos 1, correct {-}> Black), `1` (pos 2, wrong {-}> White), `5` (pos 3, wrong {-}> White). Total 1B, 2W. Matches.} \\
\texttt{        *   This looks very promising. The missing digit is `5` or `6`? Wait, if the set is `\{2, 1, 5\}`, we need a 4th digit. But the feedback says 1B + 2W = 3 correct digits. The 4th digit is incorrect.} \\
\texttt{        *   So the code contains `2`, `1`, and one of `5` or `6`. The other is wrong. And `3`, `4` are wrong.} \\
\texttt{        *   Wait, if `1` is in the code, and `2` is in the code.} \\
\texttt{        *   Let's test the hypothesis: Correct digits are `2`, `1`, and `5`. (Or `6`).} \\
\texttt{        *   If correct digits are `2, 1, 5`:} \\
\texttt{            *   Guess 1: `1` (W), `2` (W). (3,4 wrong). 0B, 2W. OK.} \\
\texttt{            *   Guess 2: `2` (B), `1` (W), `5` (W). (6 wrong). 1B, 2W. OK.} \\
\texttt{        *   If correct digits are `2, 1, 6`:} \\
\texttt{            *   Guess 1: `1` (W), `2` (W). (3,4 wrong). 0B, 2W. OK.} \\
\texttt{            *   Guess 2: `2` (B), `1` (W), `6` (W). (5 wrong). 1B, 2W. OK.} \\
\texttt{        *   If correct digits are `2, 5, 6` (and 1 is wrong):} \\
\texttt{            *   Guess 1: `1` (wrong), `2` (W). Only 1 White. **Contradiction**. Guess 1 had 2 Whites. So `1` MUST be in the code.} \\
\texttt{    *   Therefore, the correct digits are definitely **\{1, 2, and one of 5 or 6\}**.} \\
\texttt{    *   We know `2` is at position 1 (Black peg in Guess 2).} \\
\texttt{    *   We know `1` is in the code but NOT at position 1 (since 2 is there) and NOT at position 2 (since in Guess 2, `1` was at pos 2 and gave a White peg, meaning it's in the code but wrong pos).} \\
\texttt{    *   So `1` is at pos 3 or 4.} \\
\texttt{    *   We need to distinguish between `5` and `6`.} \\
\texttt{    *   In Guess 2, `5` was at pos 3 and `6` at pos 4. Both gave White pegs? No, we only have 2 White pegs total in Guess 2. One is `1`, the other is either `5` or `6`.} \\
\texttt{    *   So either `5` is in the code (and `6` is not) OR `6` is in the code (and `5` is not).} \\
\texttt{    *   Let's check Guess 1 again. Digits were `1, 2, 3, 4`. Whites came from `1` and `2`. So `3` and `4` are definitely NOT in the code.} \\
\texttt{    *   So the code consists of `2`, `1`, and either `5` or `6`.} \\
\texttt{    *   Positions:} \\
\texttt{        *   Pos 1: `2` (Confirmed Black).} \\
\texttt{        *   Pos 2: Not `1` (White in Guess 2). Not `2` (used).} \\
\texttt{        *   Pos 3/4: Contains `1` and the mystery digit (`5` or `6`).} \\
\\ 
\texttt{    Let's try to confirm the mystery digit. I'll guess a configuration that puts `1` in a new spot and keeps `2` at pos 1, while testing `5` vs `6`.} \\
\texttt{    Actually, I can just try to place `1` and `5` in the remaining spots.} \\
\texttt{    Possible codes:} \\
\texttt{    1. `2 3 1 5` (3 is wrong, so this fails) {-}> `2 X 1 5` where X is not 1,2,3,4,5,6? No, digits 1{-}6 no dupes.} \\
\texttt{    Available digits for the code: \{1, 2, 5, 6\} minus one.} \\
\texttt{    If the set is \{1, 2, 5\}, the code is a permutation of 1, 2, 5 plus a wrong digit? No, 4 digits total.} \\
\texttt{    Wait, "1B + 2W" means 3 correct digits. The 4th digit is wrong.} \\
\texttt{    The wrong digit in Guess 2 was `6` (if set is 1,2,5) or `5` (if set is 1,2,6).} \\
\texttt{    The wrong digits in Guess 1 were `3` and `4`.} \\
\texttt{    So the code is composed of \{1, 2, X, Y\} where X, Y are from \{5, 6\} and one of them is the "wrong" one from the perspective of the guess?} \\
\texttt{    No.} \\
\texttt{    Let's re{-}read carefully.} \\
\texttt{    Guess 1: `1 2 3 4`. Feedback: 0B, 2W.} \\
\texttt{    {-}> Two of \{1,2,3,4\} are in the code. Two are not.} \\
\texttt{    Guess 2: `2 1 5 6`. Feedback: 1B, 2W.} \\
\texttt{    {-}> One of \{2,1,5,6\} is correct pos. Two are correct wrong pos. One is not in code.} \\
\\ 
\texttt{    Intersection analysis:} \\
\texttt{    Digits in both: 1, 2.} \\
\texttt{    Digits only in G1: 3, 4.} \\
\texttt{    Digits only in G2: 5, 6.} \\
\\ 
\texttt{    From G1 (0B, 2W):} \\
\texttt{    {-} If 3 or 4 were in the code, they would contribute to the 2W.} \\
\texttt{    {-} If 1 and 2 are in the code, they contribute to the 2W.} \\
\texttt{    {-} If 1 and 2 are in the code, then 3 and 4 are NOT.} \\
\texttt{    {-} If 3 and 4 are in the code, then 1 and 2 are NOT.} \\
\\ 
\texttt{    From G2 (1B, 2W):} \\
\texttt{    {-} If 5 and 6 are in the code, they contribute to the 2W.} \\
\texttt{    {-} If 1 and 2 are in the code, they contribute to the 2W (plus the Black).} \\
\\ 
\texttt{    Let's test the hypothesis: **\{1, 2\} are in the code.**} \\
\texttt{    {-} Then from G1, 3 and 4 are out.} \\
\texttt{    {-} From G2, we have 1B, 2W.} \\
\texttt{    {-} We know `2` is at pos 1 in G2. If `2` is the Black peg, then `2` is at pos 1.} \\
\texttt{    {-} Then `1` is at pos 2 in G2. Since `1` is in the code (from hypothesis), and it's not at pos 1 (occupied by 2), it must be a White peg. That's 1 White.} \\
\texttt{    {-} We need 1 more White peg. It must be `5` or `6`.} \\
\texttt{    {-} So the code contains \{1, 2, 5\} or \{1, 2, 6\}.} \\
\texttt{    {-} The 4th digit is the one NOT in the set \{1, 2, 5, 6\}? No, the code has 4 digits.} \\
\texttt{    {-} If the code has \{1, 2, 5\}, the 4th digit must be from \{3, 4, 6\}?} \\
\texttt{    {-} But G1 said 3 and 4 are NOT in the code (because 1 and 2 accounted for the 2 Whites).} \\
\texttt{    {-} So the 4th digit must be `6` (if the set is 1,2,5) or `5` (if the set is 1,2,6).} \\
\texttt{    {-} BUT, if the 4th digit is `6`, then in G2 (`2 1 5 6`), `6` is present.} \\
\texttt{    {-} If the code is \{1, 2, 5, 6\}, then G2 has 4 correct digits.} \\
\texttt{    {-} G2 Feedback: 1B, 2W. Total 3 correct.} \\
\texttt{    {-} So the code has exactly 3 correct digits.} \\
\texttt{    {-} This means one of \{1, 2, 5, 6\} is NOT in the code.} \\
\texttt{    {-} We established 1 and 2 are likely in.} \\
\texttt{    {-} So either 5 is out (code has 1,2,6) or 6 is out (code has 1,2,5).} \\
\texttt{    {-} And the 4th digit? The code has 4 digits.} \\
\texttt{    {-} If the code has \{1, 2, 5, X\}, and X is not 1,2,5.} \\
\texttt{    {-} From G1, 3 and 4 are out.} \\
\texttt{    {-} So X must be 6.} \\
\texttt{    {-} So the code is \{1, 2, 5, 6\}.} \\
\texttt{    {-} But if the code is \{1, 2, 5, 6\}, then G2 (`2 1 5 6`) has 4 correct digits.} \\
\texttt{    {-} Feedback for G2 is 1B, 2W = 3 pegs.} \\
\texttt{    {-} Contradiction. The code cannot have 4 correct digits if the feedback only shows 3.} \\
\\ 
\texttt{    **Re{-}evaluating the "Two of \{1,2,3,4\}" logic.**} \\
\texttt{    Maybe 3 or 4 IS in the code?} \\
\texttt{    Case A: 1 and 2 are NOT in the code.} \\
\texttt{    {-} Then G1 (0B, 2W) means 3 and 4 are in the code.} \\
\texttt{    {-} G2 (1B, 2W) involves \{2, 1, 5, 6\}. Since 1 and 2 are out, the correct digits must be from \{5, 6\}.} \\
\texttt{    {-} Max correct digits from \{5, 6\} is 2.} \\
\texttt{    {-} Feedback is 1B + 2W = 3 correct.} \\
\texttt{    {-} Impossible to get 3 correct digits from a set of 2.} \\
\texttt{    {-} **Conclusion:** 1 and 2 MUST be in the code.} \\
\\ 
\texttt{    So, 1 and 2 are in the code.} \\
\texttt{    From G1 (0B, 2W): Since 1 and 2 are in, and we have 2 Whites, 3 and 4 are OUT.} \\
\texttt{    From G2 (1B, 2W):} \\
\texttt{    {-} Correct digits count = 3.} \\
\texttt{    {-} We have 1 and 2. We need 1 more.} \\
\texttt{    {-} The candidates are 5 or 6.} \\
\texttt{    {-} So the code contains \{1, 2, 5\} OR \{1, 2, 6\}.} \\
\texttt{    {-} The code has 4 digits. The 4th digit must be one of the remaining numbers.} \\
\texttt{    {-} Remaining numbers: \{3, 4, 6\} (if 5 is in) or \{3, 4, 5\} (if 6 is in).} \\
\texttt{    {-} But we know 3 and 4 are OUT from G1.} \\
\texttt{    {-} So the 4th digit must be the one NOT chosen from \{5, 6\}.} \\
\texttt{    {-} If the code has \{1, 2, 5\}, the 4th digit must be 6.} \\
\texttt{    {-} If the code has \{1, 2, 6\}, the 4th digit must be 5.} \\
\texttt{    {-} This leads back to the code being \{1, 2, 5, 6\}.} \\
\texttt{    {-} If the code is \{1, 2, 5, 6\}, then G2 (`2 1 5 6`) has 4 correct digits.} \\
\texttt{    {-} Why did G2 give 1B, 2W (3 pegs)?} \\
\texttt{    {-} Is it possible I misinterpreted the feedback? "1 black peg(s), 2 white peg(s)". Sum = 3.} \\
\texttt{    {-} Is it possible one of the digits is duplicated in the code? "no duplicates".} \\
\texttt{    {-} Is it possible my deduction that "1 and 2 are in the code" is wrong?} \\
\texttt{    {-} Let's re{-}run Case A: 1 and 2 are NOT in the code.} \\
\texttt{        {-} G1: 0B, 2W. Correct digits from \{3, 4\}. So 3 and 4 are in the code.} \\
\texttt{        {-} G2: 1B, 2W. Digits \{2, 1, 5, 6\}. 1 and 2 are out. So correct digits must be from \{5, 6\}.} \\
\texttt{        {-} Max correct from \{5, 6\} is 2.} \\
\texttt{        {-} Feedback requires 3 correct.} \\
\texttt{        {-} Contradiction.} \\
\texttt{    {-} So 1 and 2 MUST be in the code.} \\
\\ 
\texttt{    Is it possible that 3 or 4 is in the code AND 1 or 2 is in the code?} \\
\texttt{    {-} G1: 0B, 2W. Total 2 correct.} \\
\texttt{    {-} If 1 and 2 are in, that's 2 correct. So 3 and 4 must be out.} \\
\texttt{    {-} If 3 and 4 are in, that's 2 correct. So 1 and 2 must be out.} \\
\texttt{    {-} We proved 1 and 2 must be in. So 3 and 4 are out.} \\
\\ 
\texttt{    So the set of correct digits is \{1, 2, X\}.} \\
\texttt{    G2 has 3 correct digits. So X is either 5 or 6.} \\
\texttt{    The code has 4 digits. Where does the 4th come from?} \\
\texttt{    The only available digits are 1, 2, 3, 4, 5, 6.} \\
\texttt{    3, 4 are out.} \\
\texttt{    1, 2 are in.} \\
\texttt{    X is in (5 or 6).} \\
\texttt{    The 4th digit must be the other of 5 or 6.} \\
\texttt{    So the code IS \{1, 2, 5, 6\}.} \\
\\ 
\texttt{    **Why the discrepancy in G2?**} \\
\texttt{    Code: \{1, 2, 5, 6\}.} \\
\texttt{    Guess: {[}2 1 5 6{]}.} \\
\texttt{    {-} Pos 1: Guess 2, Code 2? (Maybe 2 is at pos 1). {-}> Black.} \\
\texttt{    {-} Pos 2: Guess 1, Code 1? (Maybe 1 is at pos 3 or 4). {-}> White.} \\
\texttt{    {-} Pos 3: Guess 5, Code 5? (Maybe 5 is at pos 4). {-}> White.} \\
\texttt{    {-} Pos 4: Guess 6, Code 6? (Maybe 6 is at pos 2). {-}> White.} \\
\texttt{    {-} Total: 1B, 3W.} \\
\texttt{    {-} But feedback was 1B, 2W.} \\
\\ 
\texttt{    This implies that in the guess {[}2 1 5 6{]}, only 3 digits matched the code.} \\
\texttt{    This means the code does NOT contain all of \{1, 2, 5, 6\}.} \\
\texttt{    It contains exactly 3 of them.} \\
\texttt{    But the code length is 4.} \\
\texttt{    So the code contains 3 digits from \{1, 2, 5, 6\} and 1 digit from \{3, 4\}?} \\
\texttt{    But G1 said 3 and 4 are out (because 1 and 2 accounted for the 2 Whites).} \\
\\ 
\texttt{    **Is it possible that 1 and 2 are NOT both in the code?**} \\
\texttt{    Let's reconsider G1: 0B, 2W.} \\
\texttt{    Possibilities for the 2 correct digits:} \\
\texttt{    1. \{1, 2\}} \\
\texttt{    2. \{1, 3\}} \\
\texttt{    3. \{1, 4\}} \\
\texttt{    4. \{2, 3\}} \\
\texttt{    5. \{2, 4\}} \\
\texttt{    6. \{3, 4\}} \\
\\ 
\texttt{    Now check G2: 1B, 2W. (Total 3 correct).} \\
\texttt{    Guess: {[}2 1 5 6{]}.} \\
\texttt{    {-} If the correct set is \{3, 4\}:} \\
\texttt{        {-} G2 has no 3 or 4. Max correct = 0. Fail.} \\
\texttt{    {-} If the correct set includes 1 or 2:} \\
\texttt{        {-} We need 3 correct digits in G2.} \\
\texttt{        {-} G2 digits: 2, 1, 5, 6.} \\
\texttt{        {-} If the correct set is \{1, 2, 3\}:} \\
\texttt{            {-} G2 has 1, 2. Missing 3. Max correct = 2. Fail.} \\
\texttt{        {-} If the correct set is \{1, 2, 5\}:} \\
\texttt{            {-} G2 has 1, 2, 5. Total 3. Possible.} \\
\texttt{        {-} If the correct set is \{1, 2, 6\}:} \\
\texttt{            {-} G2 has 1, 2, 6. Total 3. Possible.} \\
\texttt{        {-} If the correct set is \{1, 3, 5\}:} \\
\texttt{            {-} G2 has 1, 5. Missing 3. Max correct = 2. Fail.} \\
\texttt{        {-} If the correct set is \{2, 3, 5\}:} \\
\texttt{            {-} G2 has 2, 5. Missing 3. Max correct = 2. Fail.} \\
\\ 
\texttt{    So the ONLY} \\
            }
        }
    }
    & & \\ \\

    \theutterance \stepcounter{utterance}  
    & & & \multicolumn{4}{p{0.6\linewidth}}{
        \cellcolor[rgb]{0.9,0.9,0.9}{
            \makecell[{{p{\linewidth}}}]{
                \texttt{\tiny{[P1$\langle$GM]}}
                \texttt{{[}GAME{]} You attempted an invalid move. Reason: You have already guessed {[}2, 1, 5, 6{]}. Please try a different guess. Please resubmit a valid move and remember to follow the game rules to avoid penalties.} \\
            }
        }
    }
     \\ \\

    \theutterance \stepcounter{utterance}  
    & \multicolumn{4}{p{0.6\linewidth}}{
        \cellcolor[rgb]{0.561,0.737,0.561}{
            \makecell[{{p{\linewidth}}}]{
                \texttt{\tiny{[P1$\rangle$GM]}}
                \texttt{Based on the previous feedback:} \\
\texttt{1.  **Guess 1 `{[}1 2 3 4{]}`**: 0 Black, 2 White. This means exactly two digits from \{1, 2, 3, 4\} are in the code, and they are in the wrong positions. The other two are not in the code.} \\
\texttt{2.  **Guess 2 `{[}2 1 5 6{]}`**: 1 Black, 2 White. This means exactly three digits from \{2, 1, 5, 6\} are in the code. One is in the correct position (Black), and two are in the wrong positions (White).} \\
\\ 
\texttt{**Deduction:**} \\
\texttt{*   From Guess 1, we know the code contains 2 digits from \{1, 2, 3, 4\}.} \\
\texttt{*   From Guess 2, we know the code contains 3 digits from \{1, 2, 5, 6\}.} \\
\texttt{*   The intersection of these sets is \{1, 2\}.} \\
\texttt{*   If the code contained 3 or 4, they would have contributed to the count in Guess 2, but Guess 2 only has digits \{1, 2, 5, 6\}. Since Guess 2 yielded 3 correct digits, and it shares \{1, 2\} with Guess 1, the most logical overlap is that **1 and 2 are in the code**.} \\
\texttt{*   If 1 and 2 are in the code, then from Guess 1 (which had 2 Whites), the other two digits (3 and 4) must be **incorrect**.} \\
\texttt{*   So, the code consists of \{1, 2\} plus two digits from \{5, 6\}.} \\
\texttt{*   From Guess 2 (1 Black, 2 White), we have 3 correct digits total. Since we know 1 and 2 are correct, we need exactly one more correct digit from \{5, 6\}.} \\
\texttt{*   Therefore, the code contains \{1, 2\} and **one** of \{5, 6\}. The fourth digit must be the one from \{5, 6\} that is *not* in the code? Wait, the code has 4 digits.} \\
\texttt{    *   Correct digits identified: 1, 2, and (5 or 6). That's 3 digits.} \\
\texttt{    *   We need a 4th digit.} \\
\texttt{    *   Available digits: 1, 2, 3, 4, 5, 6.} \\
\texttt{    *   Eliminated: 3, 4 (from Guess 1 logic).} \\
\texttt{    *   So the 4th digit must be the remaining one from \{5, 6\}.} \\
\texttt{    *   This implies the code is a permutation of \{1, 2, 5, 6\}.} \\
\texttt{    *   Let's re{-}verify the "3 correct digits" logic. If the code is \{1, 2, 5, 6\}, then Guess 2 `{[}2 1 5 6{]}` should have 4 correct digits. But the feedback was 1B + 2W = 3 pegs.} \\
\texttt{    *   This implies one of the digits in \{1, 2, 5, 6\} is **NOT** in the code.} \\
\texttt{    *   But we established 1 and 2 MUST be in the code.} \\
\texttt{    *   So either 5 is out (and 6 is in) OR 6 is out (and 5 is in).} \\
\texttt{    *   If 5 is out, the code has \{1, 2, 6, X\}. X must be from \{3, 4\}. But 3 and 4 were eliminated by Guess 1.} \\
\texttt{    *   **Contradiction?** Let's re{-}read Guess 1. "0 Black, 2 White".} \\
\texttt{        *   Maybe the correct digits are \{3, 4\}?} \\
\texttt{        *   If correct are \{3, 4\}, then 1 and 2 are out.} \\
\texttt{        *   Then Guess 2 `{[}2 1 5 6{]}` has 0 correct digits from \{1, 2\}. It would need 3 correct from \{5, 6\}. Impossible (only 2 digits).} \\
\texttt{        *   So \{1, 2\} MUST be in the code.} \\
\texttt{    *   Maybe the correct digits in Guess 1 are \{1, 3\}?} \\
\texttt{        *   Then 2, 4 are out.} \\
\texttt{        *   Guess 2 `{[}2 1 5 6{]}`: 1 is in. 2 is out. We need 2 more correct from \{5, 6\}. So \{5, 6\} must be in.} \\
\texttt{        *   Code digits: \{1, 3, 5, 6\}.} \\
\texttt{        *   Check Guess 1 `{[}1 2 3 4{]}`:} \\
\texttt{            *   1 is in (pos 1). 3 is in (pos 3).} \\
\texttt{            *   If code is \{1, 3, 5, 6\}, where are they?} \\
\texttt{            *   Guess 1 has 1 at pos 1, 3 at pos 3.} \\
\texttt{            *   If 1 is at pos 1 in code {-}> Black. But feedback is 0 Black. So 1 is NOT at pos 1.} \\
\texttt{            *   If 3 is at pos 3 in code {-}> Black. Feedback 0 Black. So 3 is NOT at pos 3.} \\
\texttt{            *   So 1 and 3 are in the code but in wrong positions. That gives 2 Whites. This works!} \\
\texttt{        *   Check Guess 2 `{[}2 1 5 6{]}`:} \\
\texttt{            *   Digits: 2 (out), 1 (in), 5 (in), 6 (in).} \\
\texttt{            *   We have 3 correct digits: 1, 5, 6.} \\
\texttt{            *   Feedback: 1 Black, 2 White.} \\
\texttt{            *   Positions in Guess 2:} \\
\texttt{                *   Pos 1: 2 (Wrong).} \\
\texttt{                *   Pos 2: 1. Is 1 at pos 2 in code? If yes {-}> Black.} \\
\texttt{                *   Pos 3: 5. Is 5 at pos 3?} \\
\texttt{                *   Pos 4: 6. Is 6 at pos 4?} \\
\texttt{            *   We need exactly 1 Black and 2 Whites among \{1, 5, 6\}.} \\
\texttt{            *   This scenario (\{1, 3, 5, 6\}) is possible.} \\
\\ 
\texttt{    *   Let's try another combination for Guess 1: \{2, 3\}?} \\
\texttt{        *   Then 1, 4 are out.} \\
\texttt{        *   Guess 2 `{[}2 1 5 6{]}`: 2 is in. 1 is out. Need 2 more from \{5, 6\}. So \{5, 6\} in.} \\
\texttt{        *   Code: \{2, 3, 5, 6\}.} \\
\texttt{        *   Guess 1 `{[}1 2 3 4{]}`: 2 at pos 2, 3 at pos 3. Both wrong pos (0 Black). 2 Whites. Works.} \\
\texttt{        *   Guess 2 `{[}2 1 5 6{]}`: 2 at pos 1, 5 at pos 3, 6 at pos 4.} \\
\texttt{        *   Need 1 Black, 2 Whites from \{2, 5, 6\}.} \\
\texttt{        *   This is also possible.} \\
\\ 
\texttt{    *   Let's try \{1, 2\} again.} \\
\texttt{        *   Code: \{1, 2, 5, 6\}.} \\
\texttt{        *   Guess 1 `{[}1 2 3 4{]}`: 1 at pos 1, 2 at pos 2. Both wrong pos. 2 Whites. Works.} \\
\texttt{        *   Guess 2 `{[}2 1 5 6{]}`: 2 at pos 1, 1 at pos 2, 5 at pos 3, 6 at pos 4.} \\
\texttt{        *   All 4 digits are in the code.} \\
\texttt{        *   Feedback must be 1B + 3W = 4 pegs.} \\
\texttt{        *   But feedback was 1B + 2W = 3 pegs.} \\
\texttt{        *   **Conclusion**: The code does NOT contain all of \{1, 2, 5, 6\}.} \\
\texttt{        *   Therefore, the set of correct digits is NOT \{1, 2, 5, 6\}.} \\
\texttt{        *   This eliminates the case where 1 and 2 are the pair from Guess 1.} \\
\texttt{        *   So the pair from Guess 1 must include a digit from \{3, 4\}.} \\
\texttt{        *   Possibilities: \{1, 3\}, \{1, 4\}, \{2, 3\}, \{2, 4\}.} \\
\texttt{        *   As analyzed:} \\
\texttt{            *   If \{1, 3\}: Code has \{1, 3, 5, 6\}.} \\
\texttt{            *   If \{2, 3\}: Code has \{2, 3, 5, 6\}.} \\
\texttt{            *   If \{1, 4\}: Code has \{1, 4, 5, 6\}.} \\
\texttt{            *   If \{2, 4\}: Code has \{2, 4, 5, 6\}.} \\
\\ 
\texttt{    *   Let's test **\{1, 3, 5, 6\}**:} \\
\texttt{        *   Guess 2 `{[}2 1 5 6{]}`:} \\
\texttt{            *   2 is wrong.} \\
\texttt{            *   1, 5, 6 are correct.} \\
\texttt{            *   We need 1 Black, 2 Whites.} \\
\texttt{            *   If 1 is at pos 2 (Black), then 5 and 6 are wrong pos (White).} \\
\texttt{            *   Code could be `? 1 ? ?` with 5, 6 in pos 3, 4 swapped? Or 5 at 4, 6 at 3?} \\
\texttt{            *   If code is `3 1 6 5`:} \\
\texttt{                *   G1 `{[}1 2 3 4{]}`: 1 (pos 1 vs 2), 3 (pos 3 vs 1). 2 Whites. OK.} \\
\texttt{                *   G2 `{[}2 1 5 6{]}`: 1 (pos 2 match {-}> B), 5 (pos 3 vs 4 {-}> W), 6 (pos 4 vs 3 {-}> W). 1B, 2W. OK.} \\
\texttt{            *   This is a strong candidate.} \\
\\ 
\texttt{    *   Let's test **\{2, 3, 5, 6\}**:} \\
\texttt{        *   Guess 2 `{[}2 1 5 6{]}`:} \\
\texttt{            *   2, 5, 6 correct. 1 wrong.} \\
\texttt{            *   Need 1B, 2W.} \\
\texttt{            *   If 2 is at pos 1 (Black), then 5, 6 are White.} \\
\texttt{            *   Code could be `2 3 6 5`.} \\
\texttt{            *   Check G1 `{[}1 2 3 4{]}`: 2 (pos 2 vs 1), 3 (pos 3 vs 2). 2 Whites. OK.} \\
\texttt{            *   Check G2 `{[}2 1 5 6{]}`: 2 (pos 1 match {-}> B), 5 (pos 3 vs 4 {-}> W), 6 (pos 4 vs 3 {-}> W). 1B, 2W. OK.} \\
\texttt{            *   Also a strong candidate.} \\
\\ 
\texttt{    *   How to distinguish between \{1, 3, 5, 6\} and \{2, 3, 5, 6\}?} \\
\texttt{        *   In \{1, 3, 5, 6\}, digit 2 is NOT in the code.} \\
\texttt{        *   In \{2, 3, 5, 6\}, digit 1 is NOT in the code.} \\
\texttt{        *   I need to guess a number that includes 1 and 2 to see which one yields a White peg.} \\
\texttt{        *   If I guess `{[}1 2 3 5{]}`:} \\
\texttt{            *   If code is \{1, 3, 5, 6\}: 1, 3, 5 are in. 2 is out.} \\
\texttt{                *   Positions: 1 (pos 1), 2 (out), 3 (pos 3), 5 (pos 4).} \\
\texttt{                *   If code is `3 1 6 5`: 1 (pos 2 vs 1 {-}> W), 3 (pos 1 vs 3 {-}> W), 5 (pos 4 match {-}> B). Total 1B, 2W.} \\
\texttt{            *   If code is \{2, 3, 5, 6\}: 2, 3, 5 are in. 1 is out.} \\
\texttt{                *   Positions: 1 (out), 2 (pos 2), 3 (pos 3), 5 (pos 4).} \\
\texttt{                *   If code is `2 3 6 5`: 2 (pos 2 vs 1 {-}> W), 3 (pos 3 match {-}> B), 5 (pos 4 match {-}> B). Total 2B, 1W.} \\
\texttt{        *   This seems complicated to predict without knowing the exact permutation.} \\
\texttt{        *   Let's try a simpler approach. I suspect the code is one of the permutations of \{1, 3, 5, 6\} or \{2, 3, 5, 6\}.} \\
\texttt{        *   Let's guess `{[}3 5 1 2{]}`.} \\
\texttt{            *   If code is \{1, 3, 5, 6\}: 3, 5, 1 are in. 2 is out.} \\
\texttt{            *   If code is \{2, 3, 5, 6\}: 3, 5, 2 are in. 1 is out.} \\
\texttt{            *   This doesn't isolate 1 vs 2 well because both guesses have 3 digits from the set.} \\
\texttt{        *   Let's try to force a position.} \\
\texttt{        *   Hypothesis: The code is `3 1 6 5` (from \{1,3,5,6\} scenario).} \\
\texttt{        *   Hypothesis: The code is `2 3 6 5` (from \{2,3,5,6\} scenario).} \\
\texttt{        *   Notice the pattern `6 5` at the end?} \\
\texttt{        *   Let's guess `{[}3 5 6 1{]}`.} \\
\texttt{            *   If \{1, 3, 5, 6\} and code is `3 1 6 5`:} \\
\texttt{                *   3 (pos 1 match {-}> B).} \\
\texttt{                *   5 (pos 2 vs 4 {-}> W).} \\
\texttt{                *   6 (pos 3 match {-}> B).} \\
\texttt{                *   1 (pos 4 vs 2 {-}> W).} \\
\texttt{                *   Result: 2B, 2W.} \\
\texttt{            *   If \{2, 3, 5, 6\} and code is `2 3 6 5`:} \\
\texttt{                *   3 (pos 1 vs 2 {-}> W).} \\
\texttt{                *   5 (pos 2 vs 4 {-}> W).} \\
\texttt{                *   6 (pos 3 match {-}> B).} \\
\texttt{                *   1 (out).} \\
\texttt{                *   Result: 1B, 2W.} \\
\texttt{        *   This is getting too speculative. Let's go with the most likely set based on standard Mastermind logic where "1 and 2" are often the pivot. But we proved \{1,2,5,6\} is impossible.} \\
\texttt{        *   Let's try to confirm if 1 is in the code.} \\
\texttt{        *   Guess: `{[}1 3 5 6{]}`.} \\
\texttt{            *   If 1 is in (Set \{1,3,5,6\}): 1, 3, 5, 6 all in. 4 correct digits.} \\
\texttt{                *   Feedback will be X Black, Y White. Sum = 4.} \\
\texttt{            *   If 1 is out (Set \{2,3,5,6\}): 3, 5, 6 in. 1 out. 3 correct digits.} \\
\texttt{                *   Feedback will be X Black, Y White. Sum = 3.} \\
\texttt{        *   This is the perfect discriminator!} \\
\texttt{        *   If Sum = 4, then 1 is in the code.} \\
\texttt{        *   If Sum = 3, then 1 is NOT in the code (so 2 is in).} \\
\\ 
\texttt{    Let's submit `{[}1 3 5 6{]}`.} \\
\\ 
\texttt{{[}1 3 5 6{]}} \\
            }
        }
    }
    & & \\ \\

    \theutterance \stepcounter{utterance}  
    & & & \multicolumn{4}{p{0.6\linewidth}}{
        \cellcolor[rgb]{0.9,0.9,0.9}{
            \makecell[{{p{\linewidth}}}]{
                \texttt{\tiny{[P1$\langle$GM]}}
                \texttt{{[}GAME{]} Submitted {[}1 3 5 6{]}. Feedback: 1 black peg(s), 1 white peg(s).} \\
            }
        }
    }
     \\ \\

    \theutterance \stepcounter{utterance}  
    & \multicolumn{4}{p{0.6\linewidth}}{
        \cellcolor[rgb]{0.561,0.737,0.561}{
            \makecell[{{p{\linewidth}}}]{
                \texttt{\tiny{[P1$\rangle$GM]}}
                \texttt{Based on the feedback from `{[}1 3 5 6{]}` (1 Black, 1 White), the total number of correct digits is 2.} \\
\texttt{This confirms my hypothesis that the code contains exactly two digits from the set \{1, 3, 5, 6\} and two digits from the set \{2, 4\} (since 3 and 4 were the only other candidates, and we know 3 is likely in based on previous logic, but let's re{-}evaluate).} \\
\\ 
\texttt{Wait, let's re{-}trace the logic carefully.} \\
\texttt{1.  **Guess 1 `{[}1 2 3 4{]}`**: 0B, 2W. (2 correct digits from \{1,2,3,4\}).} \\
\texttt{2.  **Guess 2 `{[}2 1 5 6{]}`**: 1B, 2W. (3 correct digits from \{1,2,5,6\}).} \\
\texttt{3.  **Guess 3 `{[}1 3 5 6{]}`**: 1B, 1W. (2 correct digits from \{1,3,5,6\}).} \\
\\ 
\texttt{Let's analyze the intersection of correct digits.} \\
\texttt{From G2, we know there are 3 correct digits in \{1, 2, 5, 6\}.} \\
\texttt{From G3, we know there are 2 correct digits in \{1, 3, 5, 6\}.} \\
\\ 
\texttt{Let \$S\_1 = \textbackslash{}\{1, 2, 3, 4\textbackslash{}\}\$, \$S\_2 = \textbackslash{}\{1, 2, 5, 6\textbackslash{}\}\$, \$S\_3 = \textbackslash{}\{1, 3, 5, 6\textbackslash{}\}\$.} \\
\texttt{Let \$C\$ be the set of 4 correct digits in the code.} \\
\texttt{\$|C \textbackslash{}cap S\_1| = 2\$.} \\
\texttt{\$|C \textbackslash{}cap S\_2| = 3\$.} \\
\texttt{\$|C \textbackslash{}cap S\_3| = 2\$.} \\
\\ 
\texttt{Possible digits: 1, 2, 3, 4, 5, 6.} \\
\texttt{From \$|C \textbackslash{}cap S\_2| = 3\$, the code has 3 digits from \{1, 2, 5, 6\}.} \\
\texttt{This means exactly 1 digit from \{1, 2, 5, 6\} is NOT in the code.} \\
\texttt{The missing digit must be one of \{1, 2, 5, 6\}.} \\
\\ 
\texttt{Case A: Missing digit is 1.} \\
\texttt{Then \$C \textbackslash{}cap \textbackslash{}\{1, 2, 5, 6\textbackslash{}\} = \textbackslash{}\{2, 5, 6\textbackslash{}\}\$.} \\
\texttt{So 2, 5, 6 are in the code.} \\
\texttt{Check G3: \$C \textbackslash{}cap \textbackslash{}\{1, 3, 5, 6\textbackslash{}\}\$. Since 1 is missing, the intersection is \$\textbackslash{}\{3, 5, 6\textbackslash{}\} \textbackslash{}cap C\$.} \\
\texttt{We know 5 and 6 are in \$C\$. So \$\textbackslash{}\{5, 6\textbackslash{}\} \textbackslash{}subset C\$.} \\
\texttt{For \$|C \textbackslash{}cap S\_3|\$ to be 2, the third element must be 3? No, if 5 and 6 are in, that's 2 digits. So 3 must NOT be in \$C\$.} \\
\texttt{So if 1 is missing, then 3 is missing.} \\
\texttt{But \$C\$ must have 4 digits.} \\
\texttt{If 1 and 3 are missing, \$C\$ must be \$\textbackslash{}\{2, 4, 5, 6\textbackslash{}\}\$.} \\
\texttt{Let's check G1: \$C \textbackslash{}cap \textbackslash{}\{1, 2, 3, 4\textbackslash{}\} = \textbackslash{}\{2, 4\textbackslash{}\}\$. Count = 2. Matches!} \\
\texttt{Let's check G2: \$C \textbackslash{}cap \textbackslash{}\{1, 2, 5, 6\textbackslash{}\} = \textbackslash{}\{2, 5, 6\textbackslash{}\}\$. Count = 3. Matches!} \\
\texttt{Let's check G3: \$C \textbackslash{}cap \textbackslash{}\{1, 3, 5, 6\textbackslash{}\} = \textbackslash{}\{5, 6\textbackslash{}\}\$. Count = 2. Matches!} \\
\texttt{So the set of digits is **\{2, 4, 5, 6\}**.} \\
\\ 
\texttt{Case B: Missing digit is 2.} \\
\texttt{Then \$C \textbackslash{}cap \textbackslash{}\{1, 2, 5, 6\textbackslash{}\} = \textbackslash{}\{1, 5, 6\textbackslash{}\}\$.} \\
\texttt{So 1, 5, 6 are in.} \\
\texttt{Check G3: \$C \textbackslash{}cap \textbackslash{}\{1, 3, 5, 6\textbackslash{}\}\$. We have 1, 5, 6. That's 3 digits.} \\
\texttt{But G3 feedback says 2 correct digits.} \\
\texttt{Contradiction. So 2 cannot be the missing digit.} \\
\\ 
\texttt{Case C: Missing digit is 5.} \\
\texttt{Then \$C \textbackslash{}cap \textbackslash{}\{1, 2, 5, 6\textbackslash{}\} = \textbackslash{}\{1, 2, 6\textbackslash{}\}\$.} \\
\texttt{Check G3: \$C \textbackslash{}cap \textbackslash{}\{1, 3, 5, 6\textbackslash{}\}\$. We have 1, 6. We need 1 more from \{3\}. So 3 must be in.} \\
\texttt{So \$C\$ contains \{1, 2, 3, 6\}.} \\
\texttt{Check G1: \$C \textbackslash{}cap \textbackslash{}\{1, 2, 3, 4\textbackslash{}\} = \textbackslash{}\{1, 2, 3\textbackslash{}\}\$. Count = 3.} \\
\texttt{But G1 feedback says 2 correct digits.} \\
\texttt{Contradiction.} \\
\\ 
\texttt{Case D: Missing digit is 6.} \\
\texttt{Similar to Case C, leads to contradiction in G1.} \\
\\ 
\texttt{**Conclusion:** The digits in the code are **\{2, 4, 5, 6\}**.} \\
\texttt{Now we need to determine the positions.} \\
\texttt{We know:} \\
\texttt{{-} Digits: 2, 4, 5, 6.} \\
\texttt{{-} From G2 `{[}2 1 5 6{]}` (1B, 2W):} \\
\texttt{    {-} Digits present: 2, 5, 6. (1 is not in code).} \\
\texttt{    {-} Positions: 2 at pos 1, 5 at pos 3, 6 at pos 4.} \\
\texttt{    {-} Feedback: 1 Black, 2 White.} \\
\texttt{    {-} This means exactly one of \{2, 5, 6\} is in the correct position.} \\
\texttt{    {-} Possibilities:} \\
\texttt{        1. 2 is at pos 1 (Black). Then 5 and 6 are wrong positions.} \\
\texttt{        2. 5 is at pos 3 (Black). Then 2 and 6 are wrong positions.} \\
\texttt{        3. 6 is at pos 4 (Black). Then 2 and 5 are wrong positions.} \\
\\ 
\texttt{{-} From G1 `{[}1 2 3 4{]}` (0B, 2W):} \\
\texttt{    {-} Digits present: 2, 4. (1, 3 not in code).} \\
\texttt{    {-} Positions: 2 at pos 2, 4 at pos 4.} \\
\texttt{    {-} Feedback: 0 Black, 2 White.} \\
\texttt{    {-} This means **2 is NOT at pos 2** and **4 is NOT at pos 4**.} \\
\\ 
\texttt{Let's combine these:} \\
\texttt{{-} We know 2 is not at pos 2.} \\
\texttt{{-} From G2 analysis:} \\
\texttt{    {-} If 2 is at pos 1 (Black): Consistent with "2 not at pos 2".} \\
\texttt{    {-} If 5 is at pos 3 (Black): Then 2 is not at pos 1.} \\
\texttt{    {-} If 6 is at pos 4 (Black): Then 2 is not at pos 1.} \\
\\ 
\texttt{Let's test the position of 4.} \\
\texttt{{-} 4 is in the code.} \\
\texttt{{-} From G1, 4 is not at pos 4.} \\
\texttt{{-} From G3 `{[}1 3 5 6{]}` (1B, 1W):} \\
\texttt{    {-} Digits present: 5, 6. (1, 3 not in code).} \\
\texttt{    {-} Positions: 5 at pos 2, 6 at pos 4.} \\
\texttt{    {-} Feedback: 1 Black, 1 White.} \\
\texttt{    {-} This means exactly one of \{5, 6\} is in the correct position.} \\
\texttt{    {-} Possibilities:} \\
\texttt{        1. 5 is at pos 2 (Black). Then 6 is wrong pos.} \\
\texttt{        2. 6 is at pos 4 (Black). Then 5 is wrong pos.} \\
\\ 
\texttt{Let's synthesize:} \\
\texttt{We have constraints on 2, 4, 5, 6.} \\
\texttt{Positions: 1, 2, 3, 4.} \\
\texttt{Constraints:} \\
\texttt{1.  2 \$\textbackslash{}neq\$ 2.} \\
\texttt{2.  4 \$\textbackslash{}neq\$ 4.} \\
\texttt{3.  From G2: Exactly one of (2=1, 5=3, 6=4) is true.} \\
\texttt{4.  From G3: Exactly one of (5=2, 6=4) is true.} \\
\\ 
\texttt{Let's analyze constraint 4:} \\
\texttt{{-} Option A: 5 is at pos 2. Then 6 \$\textbackslash{}neq\$ 4.} \\
\texttt{{-} Option B: 6 is at pos 4. Then 5 \$\textbackslash{}neq\$ 2.} \\
\\ 
\texttt{**Scenario A: 5 is at pos 2.**} \\
\texttt{{-} Then 6 \$\textbackslash{}neq\$ 4.} \\
\texttt{{-} Check G2 constraint: We have 5=3 (False, since 5=2).} \\
\texttt{{-} So we need exactly one of (2=1, 6=4) to be true.} \\
\texttt{{-} Since 6 \$\textbackslash{}neq\$ 4, then **2 must be at pos 1**.} \\
\texttt{{-} So far: Pos 1 = 2, Pos 2 = 5.} \\
\texttt{{-} Remaining digits: 4, 6. Remaining positions: 3, 4.} \\
\texttt{{-} We know 4 \$\textbackslash{}neq\$ 4. So 4 must be at pos 3.} \\
\texttt{{-} Then 6 must be at pos 4.} \\
\texttt{{-} Let's check if this violates anything.} \\
\texttt{    {-} Code: `2 5 4 6`.} \\
\texttt{    {-} Check G2 `{[}2 1 5 6{]}`:} \\
\texttt{        {-} 2 at 1 (Match {-}> B).} \\
\texttt{        {-} 1 (No).} \\
\texttt{        {-} 5 at 3 (Code has 4 {-}> No).} \\
\texttt{        {-} 6 at 4 (Match {-}> B).} \\
\texttt{        {-} Result: 2 Black, 0 White.} \\
\texttt{        {-} **Contradiction!** G2 was 1B, 2W.} \\
\texttt{{-} So Scenario A is invalid.} \\
\\ 
\texttt{**Scenario B: 6 is at pos 4.**} \\
\texttt{{-} Then 5 \$\textbackslash{}neq\$ 2.} \\
\texttt{{-} Check G2 constraint: We have 6=4 (True).} \\
\texttt{{-} So we need exactly one of (2=1, 5=3) to be False.} \\
\texttt{{-} Since 6=4 is True, we cannot have any other Black pegs in G2.} \\
\texttt{{-} So 2 \$\textbackslash{}neq\$ 1 AND 5 \$\textbackslash{}neq\$ 3.} \\
\texttt{{-} We know 6 is at pos 4.} \\
\texttt{{-} We know 4 \$\textbackslash{}neq\$ 4 (from G1). Since 6 is at 4, 4 is not at 4 anyway.} \\
\texttt{{-} Where is 4?} \\
\texttt{    {-} 4 cannot be at 4.} \\
\texttt{    {-} 4 cannot be at 1? (No constraint yet).} \\
\texttt{    {-} 4 cannot be at 2? (No constraint yet).} \\
\texttt{    {-} 4 cannot be at 3? (No constraint yet).} \\
\texttt{{-} Where is 2?} \\
\texttt{    {-} 2 \$\textbackslash{}neq\$ 2 (from G1).} \\
\texttt{    {-} 2 \$\textbackslash{}neq\$ 1 (from G2 deduction above).} \\
\texttt{    {-} So 2 must be at pos 3 or 4.} \\
\texttt{    {-} But pos 4 is taken by 6.} \\
\texttt{    {-} So **2 must be at pos 3**.} \\
\texttt{{-} Where is 5?} \\
\texttt{    {-} 5 \$\textbackslash{}neq\$ 2 (from G3 deduction).} \\
\texttt{    {-} 5 \$\textbackslash{}neq\$ 3 (from G2 deduction above).} \\
\texttt{    {-} So 5 must be at pos 1 or 4.} \\
\texttt{    {-} Pos 4 is taken by 6.} \\
\texttt{    {-} So **5 must be at pos 1**.} \\
\texttt{{-} Where is 4?} \\
\texttt{    {-} Remaining spot is pos 2.} \\
\texttt{    {-} So **4 must be at pos 2**.} \\
\texttt{{-} Proposed Code: `5 4 2 6`.} \\
\\ 
\texttt{Let's verify `5 4 2 6` against all clues.} \\
\texttt{1.  **Code**: `5 4 2 6`.} \\
\texttt{2.  **Guess 1 `{[}1 2 3 4{]}`**:} \\
\texttt{    {-} 1 (No), 2 (at pos 2 vs 3 {-}> W), 3 (No), 4 (at pos 4 vs 2 {-}> W).} \\
\texttt{    {-} Result: 0B, 2W. **Matches**.} \\
\texttt{3.  **Guess 2 `{[}2 1 5 6{]}`**:} \\
\texttt{    {-} 2 (at pos 1 vs 3 {-}> W), 1 (No), 5 (at pos 3 vs 1 {-}> W), 6 (at pos 4 vs 4 {-}> B).} \\
\texttt{    {-} Result: 1B, 2W. **Matches**.} \\
\texttt{4.  **Guess 3 `{[}1 3 5 6{]}`**:} \\
\texttt{    {-} 1 (No), 3 (No), 5 (at pos 2 vs 1 {-}> W), 6 (at pos 4 vs 4 {-}> B).} \\
\texttt{    {-} Result: 1B, 1W. **Matches**.} \\
\\ 
\texttt{The code is definitely **5 4 2 6**.} \\
\\ 
\texttt{{[}5 4 2 6{]}} \\
            }
        }
    }
    & & \\ \\

    \theutterance \stepcounter{utterance}  
    & & & \multicolumn{4}{p{0.6\linewidth}}{
        \cellcolor[rgb]{0.9,0.9,0.9}{
            \makecell[{{p{\linewidth}}}]{
                \texttt{\tiny{[P1$\langle$GM]}}
                \texttt{{[}GAME{]} Submitted {[}5 4 2 6{]}. Feedback: 1 black peg(s), 3 white peg(s).} \\
            }
        }
    }
     \\ \\

    \theutterance \stepcounter{utterance}  
    & \multicolumn{4}{p{0.6\linewidth}}{
        \cellcolor[rgb]{0.561,0.737,0.561}{
            \makecell[{{p{\linewidth}}}]{
                \texttt{\tiny{[P1$\rangle$GM]}}
                \texttt{Based on the feedback from my guess `{[}5 4 2 6{]}` (1 Black, 3 White), I have successfully identified the code.} \\
\\ 
\texttt{**Analysis:**} \\
\texttt{{-} **1 Black Peg**: This indicates one digit is correct and in the correct position.} \\
\texttt{{-} **3 White Pegs**: This indicates three digits are correct but in the wrong positions.} \\
\texttt{{-} **Total Correct Digits**: 1 + 3 = 4. Since the code is 4 digits long, **all digits in my guess are in the code**.} \\
\texttt{{-} **The Code is a permutation of \{2, 4, 5, 6\}**.} \\
\\ 
\texttt{Let's re{-}evaluate the positions based on the previous feedback and this new one:} \\
\texttt{1.  **From Guess 1 `{[}1 2 3 4{]}` (0B, 2W)**:} \\
\texttt{    {-} Digits 2 and 4 were in the code (in wrong positions).} \\
\texttt{    {-} My guess had 2 at pos 2 and 4 at pos 4.} \\
\texttt{    {-} Since they are White pegs, **2 is NOT at pos 2** and **4 is NOT at pos 4**.} \\
\\ 
\texttt{2.  **From Guess 2 `{[}2 1 5 6{]}` (1B, 2W)**:} \\
\texttt{    {-} Digits 2, 5, 6 were in the code.} \\
\texttt{    {-} My guess had 2 at pos 1, 5 at pos 3, 6 at pos 4.} \\
\texttt{    {-} Feedback: 1 Black, 2 White.} \\
\texttt{    {-} This means exactly **one** of these positions is correct:} \\
\texttt{        {-} Either **2 is at pos 1**.} \\
\texttt{        {-} Or **5 is at pos 3**.} \\
\texttt{        {-} Or **6 is at pos 4**.} \\
\\ 
\texttt{3.  **From Guess 3 `{[}1 3 5 6{]}` (1B, 1W)**:} \\
\texttt{    {-} Digits 5 and 6 were in the code.} \\
\texttt{    {-} My guess had 5 at pos 2, 6 at pos 4.} \\
\texttt{    {-} Feedback: 1 Black, 1 White.} \\
\texttt{    {-} This means exactly **one** of these positions is correct:} \\
\texttt{        {-} Either **5 is at pos 2**.} \\
\texttt{        {-} Or **6 is at pos 4**.} \\
\\ 
\texttt{4.  **From Guess 4 `{[}5 4 2 6{]}` (1B, 3W)**:} \\
\texttt{    {-} My guess: 5(pos1), 4(pos2), 2(pos3), 6(pos4).} \\
\texttt{    {-} Since all digits are correct, the 1 Black peg tells us which one is in the right spot.} \\
\texttt{    {-} The 3 White pegs tell us the other three are in the wrong spots.} \\
\\ 
\texttt{Let's test the possibilities for the Black peg in Guess 4:} \\
\texttt{{-} **Hypothesis A: 5 is the Black peg (Pos 1 is correct).**} \\
\texttt{    {-} Code starts with 5: `5 \_ \_ \_`.} \\
\texttt{    {-} Then 4, 2, 6 are White (wrong positions).} \\
\texttt{    {-} So 4 \$\textbackslash{}neq\$ 2, 2 \$\textbackslash{}neq\$ 3, 6 \$\textbackslash{}neq\$ 4.} \\
\texttt{    {-} Check Guess 3 `{[}1 3 5 6{]}`: 5 is at pos 2. If code is `5...`, then 5 is NOT at pos 2. So 5 is not the Black peg in Guess 3.} \\
\texttt{    {-} In Guess 3, we need 1 Black. Since 5 is not it, **6 must be the Black peg**.} \\
\texttt{    {-} So 6 is at pos 4.} \\
\texttt{    {-} But in Hypothesis A, we said 6 is a White peg in Guess 4 (meaning 6 \$\textbackslash{}neq\$ 4).} \\
\texttt{    {-} Contradiction: 6 cannot be at pos 4 if it's a White peg in Guess 4, but we deduced it must be at pos 4 from Guess 3.} \\
\texttt{    {-} **Hypothesis A is false.**} \\
\\ 
\texttt{{-} **Hypothesis B: 4 is the Black peg (Pos 2 is correct).**} \\
\texttt{    {-} Code: `\_ 4 \_ \_`.} \\
\texttt{    {-} Then 5, 2, 6 are White (wrong positions).} \\
\texttt{    {-} So 5 \$\textbackslash{}neq\$ 1, 2 \$\textbackslash{}neq\$ 3, 6 \$\textbackslash{}neq\$ 4.} \\
\texttt{    {-} Check Guess 3 `{[}1 3 5 6{]}`: 5 is at pos 2. Code has 4 at pos 2, so 5 \$\textbackslash{}neq\$ 2.} \\
\texttt{    {-} We need 1 Black in Guess 3. Since 5 is not it, **6 must be the Black peg**.} \\
\texttt{    {-} So 6 is at pos 4.} \\
\texttt{    {-} Check consistency with Guess 4: We said 6 \$\textbackslash{}neq\$ 4. But here 6 is at 4.} \\
\texttt{    {-} Contradiction.} \\
\texttt{    {-} **Hypothesis B is false.**} \\
\\ 
\texttt{{-} **Hypothesis C: 2 is the Black peg (Pos 3 is correct).**} \\
\texttt{    {-} Code: `\_ \_ 2 \_`.} \\
\texttt{    {-} Then 5, 4, 6 are White (wrong positions).} \\
\texttt{    {-} So 5 \$\textbackslash{}neq\$ 1, 4 \$\textbackslash{}neq\$ 2, 6 \$\textbackslash{}neq\$ 4.} \\
\texttt{    {-} Check Guess 3 `{[}1 3 5 6{]}`: 5 is at pos 2. Code has 2 at pos 3.} \\
\texttt{    {-} We need 1 Black in Guess 3.} \\
\texttt{    {-} Is 5 at pos 2? Maybe.} \\
\texttt{    {-} Is 6 at pos 4? We know 6 \$\textbackslash{}neq\$ 4 from Guess 4 analysis. So 6 is not the Black peg in Guess 3.} \\
\texttt{    {-} Therefore, **5 must be the Black peg in Guess 3**.} \\
\texttt{    {-} So 5 is at pos 2.} \\
\texttt{    {-} Current Code: `\_ 5 2 \_`.} \\
\texttt{    {-} Remaining digits: 4, 6. Remaining spots: 1, 4.} \\
\texttt{    {-} We know 6 \$\textbackslash{}neq\$ 4 (from Guess 4 White peg). So 6 must be at pos 1.} \\
\texttt{    {-} Then 4 must be at pos 4.} \\
\texttt{    {-} Proposed Code: `6 5 2 4`.} \\
\texttt{    {-} Let's verify constraints:} \\
\texttt{        {-} 4 \$\textbackslash{}neq\$ 2? Yes (4 is at 4). Wait, constraint was 4 \$\textbackslash{}neq\$ 2 (from Guess 4 White peg). Yes, 4 is at 4, not 2. OK.} \\
\texttt{        {-} 2 \$\textbackslash{}neq\$ 3? Yes, 2 is at 3. Wait, if 2 is at 3, it's a Black peg in Guess 4. That's consistent.} \\
\texttt{        {-} 6 \$\textbackslash{}neq\$ 4? Yes, 6 is at 1. OK.} \\
\texttt{    {-} Let's re{-}verify Guess 3 with `6 5 2 4`:} \\
\texttt{        {-} Guess: `1 3 5 6`.} \\
\texttt{        {-} Code: `6 5 2 4`.} \\
\texttt{        {-} 5 is at pos 2 (Match {-}> Black).} \\
\texttt{        {-} 6 is at pos 4 (Code has 4 {-}> White).} \\
\texttt{        {-} Result: 1B, 1W. **Matches**.} \\
\texttt{    {-} Let's re{-}verify Guess 2 with `6 5 2 4`:} \\
\texttt{        {-} Guess: `2 1 5 6`.} \\
\texttt{        {-} Code: `6 5 2 4`.} \\
\texttt{        {-} 2 is at pos 3 (Code has 6 {-}> White).} \\
\texttt{        {-} 5 is at pos 3 (Code has 2 {-}> White).} \\
\texttt{        {-} 6 is at pos 1 (Code has 6 {-}> Black).} \\
\texttt{        {-} Result: 1B, 2W. **Matches**.} \\
\texttt{    {-} Let's re{-}verify Guess 1 with `6 5 2 4`:} \\
\texttt{        {-} Guess: `1 2 3 4`.} \\
\texttt{        {-} Code: `6 5 2 4`.} \\
\texttt{        {-} 2 is at pos 3 (Code has 2 {-}> White).} \\
\texttt{        {-} 4 is at pos 4 (Code has 4 {-}> White).} \\
\texttt{        {-} Result: 0B, 2W. **Matches**.} \\
\\ 
\texttt{The code is **6 5 2 4**.} \\
\\ 
\texttt{{[}6 5 2 4{]}} \\
            }
        }
    }
    & & \\ \\

    \theutterance \stepcounter{utterance}  
    & & & \multicolumn{4}{p{0.6\linewidth}}{
        \cellcolor[rgb]{0.9,0.9,0.9}{
            \makecell[{{p{\linewidth}}}]{
                \texttt{\tiny{[P1$\langle$GM]}}
                \texttt{{[}GAME{]} Submitted {[}6 5 2 4{]}. Feedback: 1 black peg(s), 3 white peg(s).} \\
            }
        }
    }
     \\ \\

    \theutterance \stepcounter{utterance}  
    & \multicolumn{4}{p{0.6\linewidth}}{
        \cellcolor[rgb]{0.561,0.737,0.561}{
            \makecell[{{p{\linewidth}}}]{
                \texttt{\tiny{[P1$\rangle$GM]}}
                \texttt{Based on the feedback from my guess `{[}6 5 2 4{]}` (1 Black, 3 White), I have confirmed that the set of digits is indeed \{2, 4, 5, 6\}, but my previous deduction of the specific positions was incorrect.} \\
\\ 
\texttt{Let's re{-}evaluate the constraints with the new data:} \\
\texttt{1.  **Guess 1 `{[}1 2 3 4{]}` (0B, 2W)**: Digits 2 and 4 are in the code, but **2 is not at pos 2** and **4 is not at pos 4**.} \\
\texttt{2.  **Guess 2 `{[}2 1 5 6{]}` (1B, 2W)**: Digits 2, 5, 6 are in the code. Exactly one is in the correct position.} \\
\texttt{    *   Possibilities: (2 at 1) OR (5 at 3) OR (6 at 4).} \\
\texttt{3.  **Guess 3 `{[}1 3 5 6{]}` (1B, 1W)**: Digits 5, 6 are in the code. Exactly one is in the correct position.} \\
\texttt{    *   Possibilities: (5 at 2) OR (6 at 4).} \\
\texttt{4.  **Guess 4 `{[}5 4 2 6{]}` (1B, 3W)**: All digits are correct. Exactly one is in the correct position.} \\
\texttt{    *   Possibilities: (5 at 1) OR (4 at 2) OR (2 at 3) OR (6 at 4).} \\
\texttt{5.  **Guess 5 `{[}6 5 2 4{]}` (1B, 3W)**: All digits are correct. Exactly one is in the correct position.} \\
\texttt{    *   Possibilities: (6 at 1) OR (5 at 2) OR (2 at 3) OR (4 at 4).} \\
\\ 
\texttt{**Logical Deduction:**} \\
\texttt{*   From Guess 1, we know **4 is not at pos 4**.} \\
\texttt{*   From Guess 5, since 4 is not at pos 4, the Black peg cannot be 4. So the Black peg in Guess 5 must be **6, 5, or 2**.} \\
\texttt{*   From Guess 4, the Black peg is **5, 4, 2, or 6**.} \\
\texttt{*   Let's look at **6 at pos 4**.} \\
\texttt{    *   If 6 is at pos 4:} \\
\texttt{        *   In Guess 2 (`...6`), 6 would be Black. Then 2 and 5 must be White (wrong pos). So 2 \$\textbackslash{}neq\$ 1, 5 \$\textbackslash{}neq\$ 3.} \\
\texttt{        *   In Guess 3 (`...6`), 6 would be Black. Then 5 must be White (wrong pos). So 5 \$\textbackslash{}neq\$ 2.} \\
\texttt{        *   In Guess 4 (`...6`), 6 would be Black. Then 5, 4, 2 must be White. So 5 \$\textbackslash{}neq\$ 1, 4 \$\textbackslash{}neq\$ 2, 2 \$\textbackslash{}neq\$ 3.} \\
\texttt{        *   In Guess 5 (`6...`), 6 would be Black. Then 5, 2, 4 must be White. So 5 \$\textbackslash{}neq\$ 2, 2 \$\textbackslash{}neq\$ 3, 4 \$\textbackslash{}neq\$ 4.} \\
\texttt{        *   This seems consistent so far. If 6 is at pos 4, where are the others?} \\
\texttt{        *   We need to place 2, 4, 5 in pos 1, 2, 3.} \\
\texttt{        *   Constraints on 2: Not at 1 (from G2), Not at 3 (from G4/G5). So **2 must be at pos 2**.} \\
\texttt{        *   But Guess 1 says **2 is not at pos 2**. **Contradiction**.} \\
\texttt{    *   Therefore, **6 is NOT at pos 4**.} \\
\\ 
\texttt{*   Since 6 is not at pos 4:} \\
\texttt{    *   From Guess 3 (`...6`), 6 is not Black. So **5 must be the Black peg**.} \\
\texttt{    *   **5 is at pos 2**.} \\
\texttt{    *   Now check Guess 2 (`...5...`): 5 is at pos 3. Since 5 is actually at pos 2, 5 is a White peg.} \\
\texttt{    *   We need 1 Black in Guess 2. It can't be 5. It must be 2 or 6.} \\
\texttt{    *   We know 6 is not at pos 4. Could 6 be at pos 1?} \\
\texttt{    *   Let's check Guess 4 (`5 4 2 6`): 5 is at pos 1. But we know 5 is at pos 2. So 5 is a White peg.} \\
\texttt{    *   We need 1 Black in Guess 4. It can't be 5. It must be 4, 2, or 6.} \\
\texttt{    *   Let's check Guess 5 (`6 5 2 4`): 5 is at pos 2. This matches our finding!} \\
\texttt{    *   If 5 is the Black peg in Guess 5, then 6, 2, 4 are White.} \\
\texttt{        *   6 \$\textbackslash{}neq\$ 1.} \\
\texttt{        *   2 \$\textbackslash{}neq\$ 3.} \\
\texttt{        *   4 \$\textbackslash{}neq\$ 4 (Already known).} \\
\texttt{    *   So in Guess 5, **5 is the Black peg**.} \\
\texttt{    *   Now back to Guess 4 (`5 4 2 6`): 5 is at pos 1 (Wrong, it's at 2). So 5 is White.} \\
\texttt{    *   We need 1 Black in Guess 4. Candidates: 4, 2, 6.} \\
\texttt{    *   Back to Guess 2 (`2 1 5 6`): 5 is at pos 3 (Wrong, it's at 2). So 5 is White.} \\
\texttt{    *   We need 1 Black in Guess 2. Candidates: 2, 6.} \\
\texttt{    *   Let's test **2 at pos 1**:} \\
\texttt{        *   If 2 is at pos 1:} \\
\texttt{            *   Guess 2: 2 is Black. (Consistent).} \\
\texttt{            *   Guess 4: 2 is at pos 3 (Wrong). 4 is at pos 2 (Wrong, 5 is there). 6 is at pos 4 (Wrong).} \\
\texttt{                *   Wait, if 2 is at 1, 5 is at 2. Remaining spots 3, 4 for 4, 6.} \\
\texttt{                *   Guess 4 has 4 at pos 2 (Wrong), 2 at pos 3 (Wrong), 6 at pos 4.} \\
\texttt{                *   If 6 is at pos 4, then 6 is Black in Guess 4.} \\
\texttt{                *   But we established 6 is not at pos 4 (led to contradiction earlier).} \\
\texttt{                *   So 6 is not at pos 4. Thus 6 must be at pos 3.} \\
\texttt{                *   Then 4 must be at pos 4.} \\
\texttt{                *   But Guess 1 says **4 is not at pos 4**. **Contradiction**.} \\
\texttt{    *   Therefore, **2 is NOT at pos 1**.} \\
\texttt{    *   So in Guess 2, the Black peg must be **6**.} \\
\texttt{    *   **6 is at pos 1**.} \\
\\ 
\texttt{*   Current State:} \\
\texttt{    *   Pos 1: **6**} \\
\texttt{    *   Pos 2: **5**} \\
\texttt{    *   Remaining digits: 2, 4. Remaining spots: 3, 4.} \\
\texttt{    *   Constraints:} \\
\texttt{        *   2 \$\textbackslash{}neq\$ 2 (from G1). (Satisfied, 2 is at 3 or 4).} \\
\texttt{        *   4 \$\textbackslash{}neq\$ 4 (from G1). So **4 must be at pos 3**.} \\
\texttt{        *   Then **2 must be at pos 4**.} \\
\\ 
\texttt{*   Proposed Code: **6 5 4 2**.} \\
\\ 
\texttt{Let's verify `6 5 4 2`:} \\
\texttt{1.  **G1 `{[}1 2 3 4{]}`**:} \\
\texttt{    *   2 at pos 2 (Code has 5) {-}> W.} \\
\texttt{    *   4 at pos 4 (Code has 2) {-}> W.} \\
\texttt{    *   Result: 0B, 2W. **OK**.} \\
\texttt{2.  **G2 `{[}2 1 5 6{]}`**:} \\
\texttt{    *   2 at pos 1 (Code has 6) {-}> W.} \\
\texttt{    *   5 at pos 3 (Code has 4) {-}> W.} \\
\texttt{    *   6 at pos 4 (Code has 2) {-}> W.} \\
\texttt{    *   Wait, where is the Black peg?} \\
\texttt{    *   My code is `6 5 4 2`. Guess is `2 1 5 6`.} \\
\texttt{    *   Pos 1: 2 vs 6 (No).} \\
\texttt{    *   Pos 2: 1 vs 5 (No).} \\
\texttt{    *   Pos 3: 5 vs 4 (No).} \\
\texttt{    *   Pos 4: 6 vs 2 (No).} \\
\texttt{    *   Result: 0B, 4W.} \\
\texttt{    *   **Contradiction!** G2 was 1B, 2W.} \\
\\ 
\texttt{Where did I go wrong?} \\
\texttt{Let's re{-}examine "6 is not at pos 4".} \\
\texttt{If 6 is at pos 4:} \\
\texttt{{-} G2 (`...6`): 6 is Black. 2, 5 are White.} \\
\texttt{  {-} 2 \$\textbackslash{}neq\$ 1.} \\
\texttt{  {-} 5 \$\textbackslash{}neq\$ 3.} \\
\texttt{{-} G3 (`...6`): 6 is Black. 5 is White.} \\
\texttt{  {-} 5 \$\textbackslash{}neq\$ 2.} \\
\texttt{{-} G4 (`...6`): 6 is Black. 5, 4, 2 are White.} \\
\texttt{  {-} 5 \$\textbackslash{}neq\$ 1.} \\
\texttt{  {-} 4 \$\textbackslash{}neq\$ 2.} \\
\texttt{  {-} 2 \$\textbackslash{}neq\$ 3.} \\
\texttt{{-} G5 (`6...`): 6 is Black. 5, 2, 4 are White.} \\
\texttt{  {-} 5 \$\textbackslash{}neq\$ 2.} \\
\texttt{  {-} 2 \$\textbackslash{}neq\$ 3.} \\
\texttt{  {-} 4 \$\textbackslash{}neq\$ 4.} \\
\texttt{{-} Placement:} \\
\texttt{  {-} 6 is at 4.} \\
\texttt{  {-} 2 \$\textbackslash{}neq\$ 1, 2 \$\textbackslash{}neq\$ 3. So **2 is at 2**.} \\
\texttt{  {-} But G1 says **2 \$\textbackslash{}neq\$ 2**.} \\
\texttt{  {-} Contradiction holds. 6 is definitely not at 4.} \\
\\ 
\texttt{So 5 MUST be the Black peg in G3. **5 is at pos 2**.} \\
\texttt{And 5 MUST be the Black peg in G5. **5 is at pos 2**.} \\
\texttt{This part is solid.} \\
\\ 
\texttt{Now G2 (`2 1 5 6`):} \\
\texttt{{-} 5 is at pos 3 (Wrong).} \\
\texttt{{-} We need 1 Black. Must be 2 or 6.} \\
\texttt{{-} If 2 is Black (at 1):} \\
\texttt{  {-} 2 is at 1.} \\
\texttt{  {-} G4 (`5 4 2 6`): 5 at 1 (Wrong), 4 at 2 (Wrong), 2 at 3 (Wrong), 6 at 4.} \\
\texttt{    {-} Need 1 Black. Must be 6. So 6 is at 4.} \\
\texttt{    {-} But we proved 6 cannot be at 4.} \\
\texttt{  {-} So 2 cannot be Black in G2.} \\
\texttt{{-} Therefore, **6 must be Black in G2**.} \\
\texttt{{-} **6 is at pos 1**.} \\
\\ 
\texttt{So far:} \\
\texttt{{-} Pos 1: 6} \\
\texttt{{-} Pos 2: 5} \\
\texttt{{-} Remaining: 2, 4 at pos 3, 4.} \\
\texttt{{-} G1 says 4 \$\textbackslash{}neq\$ 4. So 4 is at 3, 2 is at 4.} \\
\texttt{{-} Code: `6 5 4 2`.} \\
\texttt{{-} Why did G2 fail?} \\
\texttt{  {-} Guess: `2 1 5 6`. Code: `6 5 4 2`.} \\
\texttt{  {-} Matches: None?} \\
\texttt{  {-} 2 (pos 1) vs 6.} \\
\texttt{  {-} 1 (pos 2) vs 5.} \\
\texttt{  {-} 5 (pos 3) vs 4.} \\
\texttt{  {-} 6 (pos 4) vs 2.} \\
\texttt{  {-} 0 matches.} \\
\texttt{  {-} But G2 said 1B, 2W.} \\
\texttt{  {-} This implies my deduction that "6 is at pos 1" is wrong, OR "5 is at pos 2" is wrong.} \\
\\ 
\texttt{Let's re{-}evaluate "5 is at pos 2".} \\
\texttt{From G3 (`1 3 5 6`): 1B, 1W. Digits 5, 6.} \\
\texttt{{-} If 5 is at 2 (Black), then 6 is White (not at 4).} \\
\texttt{{-} If 6 is at 4 (Black), then 5 is White (not at 2).} \\
\texttt{From G5 (`6 5 2 4`): 1B, 3W. Digits 6, 5, 2, 4.} \\
\texttt{{-} If 6 is at 1 (Black):} \\
\texttt{  {-} Then 5, 2, 4 are White.} \\
\texttt{  {-} 5 \$\textbackslash{}neq\$ 2.} \\
\texttt{  {-} 2 \$\textbackslash{}neq\$ 3.} \\
\texttt{  {-} 4 \$\textbackslash{}neq\$ 4.} \\
\texttt{{-} If 5 is at 2 (Black):} \\
\texttt{  {-} Then 6, 2, 4 are White.} \\
\texttt{  {-} 6 \$\textbackslash{}neq\$ 1.} \\
\texttt{  {-} 2 \$\textbackslash{}neq\$ 3.} \\
\texttt{  {-} 4 \$\textbackslash{}neq\$ 4.} \\
\texttt{{-} If 2 is at 3 (Black):} \\
\texttt{  {-} Then 6, 5, 4 are White.} \\
\texttt{  {-} 6 \$\textbackslash{}neq\$ 1.} \\
\texttt{  {-} 5 \$\textbackslash{}neq\$ 2.} \\
\texttt{  {-} 4 \$\textbackslash{}neq\$ 4.} \\
\texttt{{-} If 4 is at 4 (Black):} \\
\texttt{  {-} Impossible (G1 says 4 \$\textbackslash{}neq\$ 4).} \\
\\ 
\texttt{Let's try **2 is at pos 3** (Black in G5).} \\
\texttt{{-} Code: `\_ \_ 2 \_`.} \\
\texttt{{-} G5: 2 is Black. 6, 5, 4 are White.} \\
\texttt{  {-} 6 \$\textbackslash{}neq\$ 1.} \\
\texttt{  {-} 5 \$\textbackslash{}neq\$ 2.} \\
\texttt{  {-} 4 \$\textbackslash{}neq\$ 4.} \\
\texttt{{-} G3 (`1 3 5 6`): 1B, 1W.} \\
\texttt{  {-} 5 is at 2. We know 5 \$\textbackslash{}neq\$ 2. So 5 is White.} \\
\texttt{  {-} We need 1 Black. Must be 6.} \\
\texttt{  {-} So **6 is at pos 4**.} \\
\texttt{{-} But earlier I proved 6 cannot be at 4.} \\
\texttt{  {-} Let's re{-}verify that proof.} \\
\texttt{  {-} If 6 is at 4:} \\
\texttt{    {-} G2 (`2 1 5 6`): 6 is Black. 2, 5 are White.} \\
\texttt{      {-} 2 \$\textbackslash{}neq\$ 1.} \\
\texttt{      {-} 5 \$\textbackslash{}neq\$ 3.} \\
\texttt{    {-} G1 (`1 2 3 4`): 2, 4 are White.} \\
\texttt{      {-} 2 \$\textbackslash{}neq\$ 2.} \\
\texttt{      {-} 4 \$\textbackslash{}neq\$ 4.} \\
\texttt{    {-} Current: 6 at 4.} \\
\texttt{    {-} 2 \$\textbackslash{}neq\$ 1, 2 \$\textbackslash{}neq\$ 2. So 2 is at 3. (Matches G5 hypothesis).} \\
\texttt{    {-} 5 \$\textbackslash{}neq\$ 3.} \\
\texttt{    {-} 4 \$\textbackslash{}neq\$ 4.} \\
\texttt{    {-} Spots left: 1, 2. Digits left: 4, 5.} \\
\texttt{    {-} 5 \$\textbackslash{}neq\$ 3 (irrelevant).} \\
\texttt{    {-} 4 \$\textbackslash{}neq\$ 4 (irrelevant).} \\
\texttt{    {-} We need to place 4, 5 in 1, 2.} \\
\texttt{    {-} Check G4 (`5 4 2 6`): 1B, 3W.} \\
\texttt{      {-} 6 is at 4 (Black).} \\
\texttt{      {-} 5 at 1, 4 at 2, 2 at 3.} \\
\texttt{      {-} We need 5, 4, 2 to be White.} \\
\texttt{      {-} 5 \$\textbackslash{}neq\$ 1.} \\
\texttt{      {-} 4 \$\textbackslash{}neq\$ 2.} \\
\texttt{      {-} 2 \$\textbackslash{}neq\$ 3. (We have 2 at 3). **Contradiction**.} \\
\texttt{      {-} If 2 is at 3, it is Black in G4. But G4 has only 1 Black (which is 6).} \\
\texttt{      {-} So 2 cannot be at 3 if 6 is at 4.} \\
\texttt{      {-} But 2 must be at 3 because 2 \$\textbackslash{}neq\$ 1 and 2 \$\textbackslash{}neq\$ 2.} \\
\texttt{      {-} So 6 cannot be at 4.} \\
\\ 
\texttt{Okay, so 6 is NOT at 4.} \\
\texttt{So in G3, 6 is White. So **5 must be Black in G3**.} \\
\texttt{**5 is at pos 2**.} \\
\\ 
\texttt{Now G5 (`6 5 2 4`): 1B, 3W.} \\
\texttt{{-} 5 is at 2 (Black).} \\
\texttt{{-} 6, 2, 4 are White.} \\
\texttt{{-} 6 \$\textbackslash{}neq\$ 1.} \\
\texttt{{-} 2 \$\textbackslash{}neq\$ 3.} \\
\texttt{{-} 4 \$\textbackslash{}neq\$ 4.} \\
\\ 
\texttt{Now G2 (`2 1 5 6`): 1B, 2W.} \\
\texttt{{-} 5 is at 3 (White, since 5 is at 2).} \\
\texttt{{-} Need 1 Black from 2, 6.} \\
\texttt{{-} If 2 is Black (at 1):} \\
\texttt{  {-} 2 is at 1.} \\
\texttt{  {-} G4 (`5 4 2 6`): 5 at 1 (White), 4 at 2 (White), 2 at 3 (White), 6 at 4.} \\
\texttt{    {-} Need 1 Black. Must be 6.} \\
\texttt{    {-} So 6 is at 4.} \\
\texttt{    {-} But we know 6 \$\textbackslash{}neq\$ 4. Contradiction.} \\
\texttt{{-} So **6 must be Black in G2**.} \\
\texttt{{-} **6 is at pos 1**.} \\
\\ 
\texttt{So:} \\
\texttt{{-} Pos 1: 6} \\
\texttt{{-} Pos 2: 5} \\
\texttt{{-} 2 \$\textbackslash{}neq\$ 3 (from G5).} \\
\texttt{{-} 4 \$\textbackslash{}neq\$ 4 (from G1/G5).} \\
\texttt{{-} Remaining spots: 3, 4. Digits: 2, 4.} \\
\texttt{{-} 2 cannot be at 3. So **2 is at 4**.} \\
\texttt{{-} 4 cannot be at 4. So **4 is at 3**.} \\
\texttt{{-} Code: **6 5 4 2**.} \\
\\ 
\texttt{Why did G2 fail for `6 5 4 2`?} \\
\texttt{Guess: `2 1 5 6`. Code: `6 5 4 2`.} \\
\texttt{{-} 2 (pos 1) vs 6.} \\
\texttt{{-} 1 (pos 2) vs 5.} \\
\texttt{{-} 5 (pos 3) vs 4.} \\
\texttt{{-} 6 (pos 4) vs 2.} \\
\texttt{{-} 0 matches.} \\
\texttt{{-} Feedback was 1B, 2W.} \\
\texttt{{-} This implies my deduction of positions is wrong.} \\
\texttt{{-} Is it possible 5 is NOT at 2?} \\
\texttt{{-} If 5 is not at 2, then in G3, 6 must be at 4.} \\
\texttt{{-} But we proved 6 cannot be at 4.} \\
\texttt{{-} Is it possible 6 IS at 4 and my proof was flawed?} \\
\texttt{  {-} Proof recap: If 6 at 4 {-}> 2 at 3 {-}> G4 has 2 at 3 (Black) and 6 at 4 (Black) {-}> 2 Blacks. G4 has 1 Black. Contradiction.} \\
\texttt{  {-} This seems solid.} \\
\\ 
\texttt{Is it possible the set of digits is wrong?} \\
\texttt{{-} G1: 2 correct from \{1,2,3,4\}.} \\
\texttt{{-} G2: 3 correct from \{1,2,5,6\}.} \\
\texttt{{-} G3: 2 correct from \{1,3,5,6\}.} \\
\texttt{{-} G4: 4 correct (all in guess).} \\
\texttt{{-} G5: 4 correct (all in guess).} \\
\texttt{{-} Digits in G4/G5: 2, 4, 5, 6.} \\
\texttt{{-} G1 check: \{2, 4\} in \{1,2,3,4\}. Correct.} \\
\texttt{{-} G2 check: \{2, 5, 6\} in \{1,2,5,6\}. Correct.} \\
\texttt{{-} G3 check: \{5, 6\} in \{1,3,5,6\}. Correct.} \\
\texttt{{-} The set is definitely \{2, 4, 5, 6\}.} \\
\\ 
\texttt{Let's re{-}read G2 feedback: "1 black peg(s), 2 white peg(s)".} \\
\texttt{Maybe I am misinterpreting the positions in G2?} \\
\texttt{Guess: `2 1 5 6`.} \\
\texttt{If Code is `6 5 4 2`:} \\
\texttt{{-} 2 (pos 1) != 6.} \\
\texttt{{-} 1 (pos 2) != 5.} \\
\texttt{{-} 5 (pos 3) != 4.} \\
\texttt{{-} 6 (pos 4) != 2.} \\
\texttt{{-} 0 matches.} \\
\texttt{This is a hard contradiction.} \\
\\ 
\texttt{Could the code be **4 5 2 6**?} \\
\texttt{{-} G1 (`1 2 3 4`): 4 at 4 (W), 2 at 3 (W). 0B, 2W. OK.} \\
\texttt{{-} G2 (`2 1 5 6`): 5 at 3 (B), 2 at 3 (W? No, 2 is at 3 in code, 5 is at 3 in guess).} \\
\texttt{  {-} Guess: 2, 1, 5, 6.} \\
\texttt{  {-} Code} \\
            }
        }
    }
    & & \\ \\

    \theutterance \stepcounter{utterance}  
    & & & \multicolumn{4}{p{0.6\linewidth}}{
        \cellcolor[rgb]{0.9,0.9,0.9}{
            \makecell[{{p{\linewidth}}}]{
                \texttt{\tiny{[P1$\langle$GM]}}
                \texttt{{[}GAME{]} You attempted an invalid move. Reason: You have already guessed {[}2, 1, 5, 6{]}. Please try a different guess. Please resubmit a valid move and remember to follow the game rules to avoid penalties.} \\
            }
        }
    }
     \\ \\

    \theutterance \stepcounter{utterance}  
    & \multicolumn{4}{p{0.6\linewidth}}{
        \cellcolor[rgb]{0.561,0.737,0.561}{
            \makecell[{{p{\linewidth}}}]{
                \texttt{\tiny{[P1$\rangle$GM]}}
                \texttt{Based on the feedback from my previous guess `{[}6 5 2 4{]}` (1 Black, 3 White), I know the code is a permutation of \{2, 4, 5, 6\}. The feedback indicates that exactly one digit is in the correct position, and the other three are in the wrong positions.} \\
\\ 
\texttt{Let's re{-}evaluate the constraints carefully:} \\
\texttt{1.  **Guess 1 `{[}1 2 3 4{]}` (0B, 2W)**: Digits 2 and 4 are in the code. Neither is in the correct position.} \\
\texttt{    *   \$2 \textbackslash{}neq \textbackslash{}text\{pos \} 2\$} \\
\texttt{    *   \$4 \textbackslash{}neq \textbackslash{}text\{pos \} 4\$} \\
\texttt{2.  **Guess 2 `{[}2 1 5 6{]}` (1B, 2W)**: Digits 2, 5, 6 are in the code. Exactly one is in the correct position.} \\
\texttt{    *   Possibilities: (\$2=\textbackslash{}text\{pos \} 1\$) OR (\$5=\textbackslash{}text\{pos \} 3\$) OR (\$6=\textbackslash{}text\{pos \} 4\$).} \\
\texttt{3.  **Guess 3 `{[}1 3 5 6{]}` (1B, 1W)**: Digits 5, 6 are in the code. Exactly one is in the correct position.} \\
\texttt{    *   Possibilities: (\$5=\textbackslash{}text\{pos \} 2\$) OR (\$6=\textbackslash{}text\{pos \} 4\$).} \\
\texttt{4.  **Guess 4 `{[}5 4 2 6{]}` (1B, 3W)**: All digits are in the code. Exactly one is in the correct position.} \\
\texttt{    *   Possibilities: (\$5=\textbackslash{}text\{pos \} 1\$) OR (\$4=\textbackslash{}text\{pos \} 2\$) OR (\$2=\textbackslash{}text\{pos \} 3\$) OR (\$6=\textbackslash{}text\{pos \} 4\$).} \\
\texttt{5.  **Guess 5 `{[}6 5 2 4{]}` (1B, 3W)**: All digits are in the code. Exactly one is in the correct position.} \\
\texttt{    *   Possibilities: (\$6=\textbackslash{}text\{pos \} 1\$) OR (\$5=\textbackslash{}text\{pos \} 2\$) OR (\$2=\textbackslash{}text\{pos \} 3\$) OR (\$4=\textbackslash{}text\{pos \} 4\$).} \\
\\ 
\texttt{**Deduction Steps:**} \\
\texttt{*   From Guess 1, we know \$4 \textbackslash{}neq \textbackslash{}text\{pos \} 4\$.} \\
\texttt{*   From Guess 5, since \$4 \textbackslash{}neq \textbackslash{}text\{pos \} 4\$, the Black peg cannot be 4. So the Black peg in Guess 5 is either 6, 5, or 2.} \\
\texttt{*   From Guess 3, the Black peg is either 5 (at pos 2) or 6 (at pos 4).} \\
\texttt{*   **Hypothesis: 6 is at pos 4.**} \\
\texttt{    *   If \$6 = \textbackslash{}text\{pos \} 4\$:} \\
\texttt{        *   In Guess 2, 6 is Black. Then 2 and 5 must be White (\$2 \textbackslash{}neq 1, 5 \textbackslash{}neq 3\$).} \\
\texttt{        *   In Guess 3, 6 is Black. Then 5 must be White (\$5 \textbackslash{}neq 2\$).} \\
\texttt{        *   In Guess 4, 6 is Black. Then 5, 4, 2 must be White (\$5 \textbackslash{}neq 1, 4 \textbackslash{}neq 2, 2 \textbackslash{}neq 3\$).} \\
\texttt{        *   In Guess 5, 6 is Black. Then 5, 2, 4 must be White (\$5 \textbackslash{}neq 2, 2 \textbackslash{}neq 3, 4 \textbackslash{}neq 4\$).} \\
\texttt{        *   Now place the remaining digits 2, 4, 5 in positions 1, 2, 3.} \\
\texttt{        *   We know \$2 \textbackslash{}neq 1\$ (from G2), \$2 \textbackslash{}neq 2\$ (from G1), \$2 \textbackslash{}neq 3\$ (from G4/G5).} \\
\texttt{        *   This leaves no position for 2. **Contradiction.**} \\
\texttt{    *   Therefore, **6 is NOT at pos 4**.} \\
\\ 
\texttt{*   Since \$6 \textbackslash{}neq \textbackslash{}text\{pos \} 4\$:} \\
\texttt{    *   From Guess 3, the Black peg must be **5 at pos 2**. So **\$5 = \textbackslash{}text\{pos \} 2\$**.} \\
\texttt{    *   From Guess 5, since \$5 = \textbackslash{}text\{pos \} 2\$, the Black peg is 5. This is consistent.} \\
\texttt{    *   From Guess 2, 5 is at pos 3 (Wrong, since it's at pos 2). So 5 is a White peg. We need 1 Black from 2 or 6.} \\
\texttt{        *   If 2 is Black (\$2 = \textbackslash{}text\{pos \} 1\$):} \\
\texttt{            *   In Guess 4 (`5 4 2 6`), 2 is at pos 3 (Wrong). 5 is at pos 1 (Wrong). 4 is at pos 2 (Wrong, since 5 is there). So 6 must be Black (\$6 = \textbackslash{}text\{pos \} 4\$).} \\
\texttt{            *   But we already proved \$6 \textbackslash{}neq \textbackslash{}text\{pos \} 4\$. **Contradiction.**} \\
\texttt{        *   Therefore, **6 must be the Black peg in Guess 2**. So **\$6 = \textbackslash{}text\{pos \} 1\$**.} \\
\\ 
\texttt{*   Current known positions:} \\
\texttt{    *   Pos 1: **6**} \\
\texttt{    *   Pos 2: **5**} \\
\texttt{    *   Remaining digits: 2, 4. Remaining positions: 3, 4.} \\
\texttt{    *   Constraints:} \\
\texttt{        *   \$2 \textbackslash{}neq \textbackslash{}text\{pos \} 2\$ (Satisfied).} \\
\texttt{        *   \$4 \textbackslash{}neq \textbackslash{}text\{pos \} 4\$ (From Guess 1).} \\
\texttt{        *   From Guess 5 (`6 5 2 4`), 2 is at pos 3. Since 5 is the Black peg, 2 must be a White peg. So **\$2 \textbackslash{}neq \textbackslash{}text\{pos \} 3\$**.} \\
\texttt{        *   Since \$2 \textbackslash{}neq \textbackslash{}text\{pos \} 3\$ and \$2 \textbackslash{}neq \textbackslash{}text\{pos \} 2\$ and \$2 \textbackslash{}neq \textbackslash{}text\{pos \} 1\$ (occupied by 6), **2 must be at pos 4**.} \\
\texttt{        *   Consequently, **4 must be at pos 3**.} \\
\\ 
\texttt{*   Proposed Code: **6 5 4 2**.} \\
\\ 
\texttt{Let's double check this code against all clues:} \\
\texttt{1.  **Code**: `6 5 4 2`} \\
\texttt{2.  **Guess 1 `{[}1 2 3 4{]}`**:} \\
\texttt{    *   2 is at pos 3 (Guess pos 2) {-}> White.} \\
\texttt{    *   4 is at pos 3 (Guess pos 4) {-}> White.} \\
\texttt{    *   Result: 0B, 2W. **Matches**.} \\
\texttt{3.  **Guess 2 `{[}2 1 5 6{]}`**:} \\
\texttt{    *   6 is at pos 1 (Guess pos 4) {-}> White.} \\
\texttt{    *   5 is at pos 2 (Guess pos 3) {-}> White.} \\
\texttt{    *   2 is at pos 4 (Guess pos 1) {-}> White.} \\
\texttt{    *   Wait, this gives 0B, 3W. But feedback was 1B, 2W.} \\
\texttt{    *   My deduction that "6 is the Black peg in Guess 2" implies \$6 = \textbackslash{}text\{pos \} 1\$.} \\
\texttt{    *   If \$6 = \textbackslash{}text\{pos \} 1\$, then in Guess 2 (`2 1 5 6`), 6 is at pos 4. It should be a White peg.} \\
\texttt{    *   Where is the Black peg in Guess 2?} \\
\texttt{    *   Guess 2: `2 1 5 6`. Code: `6 5 4 2`.} \\
\texttt{    *   Pos 1: 2 vs 6 (No).} \\
\texttt{    *   Pos 2: 1 vs 5 (No).} \\
\texttt{    *   Pos 3: 5 vs 4 (No).} \\
\texttt{    *   Pos 4: 6 vs 2 (No).} \\
\texttt{    *   This results in 0 matches. This contradicts the feedback "1B, 2W".} \\
\\ 
\texttt{There must be an error in my manual deduction logic regarding the "Black peg" assignment in Guess 2.} \\
\texttt{Let's re{-}examine Guess 2: `2 1 5 6` {-}> 1B, 2W.} \\
\texttt{Digits in code: 2, 5, 6.} \\
\texttt{Positions in guess: 2@1, 5@3, 6@4.} \\
\texttt{One of these is correct.} \\
\texttt{If \$6 = \textbackslash{}text\{pos \} 1\$ (from my previous deduction), then in Guess 2, 6 is at pos 4. So 6 is NOT the Black peg.} \\
\texttt{So the Black peg in Guess 2 must be 2 or 5.} \\
\texttt{If \$5 = \textbackslash{}text\{pos \} 2\$ (from Guess 3 deduction), then in Guess 2, 5 is at pos 3. So 5 is NOT the Black peg.} \\
\texttt{So the Black peg in Guess 2 MUST be **2 at pos 1**.} \\
\texttt{So **\$2 = \textbackslash{}text\{pos \} 1\$**.} \\
\\ 
\texttt{Let's restart the position assignment with **\$2 = \textbackslash{}text\{pos \} 1\$**.} \\
\texttt{*   Pos 1: **2**} \\
\texttt{*   From Guess 1 (`1 2 3 4`): 2 is at pos 2 (Wrong). Consistent.} \\
\texttt{*   From Guess 2 (`2 1 5 6`): 2 is at pos 1 (Black). Consistent.} \\
\texttt{    *   Then 5 and 6 are White.} \\
\texttt{    *   \$5 \textbackslash{}neq \textbackslash{}text\{pos \} 3\$.} \\
\texttt{    *   \$6 \textbackslash{}neq \textbackslash{}text\{pos \} 4\$.} \\
\texttt{*   From Guess 3 (`1 3 5 6`): 1B, 1W. Digits 5, 6.} \\
\texttt{    *   We need 1 Black.} \\
\texttt{    *   Can 6 be Black? If \$6 = \textbackslash{}text\{pos \} 4\$. But we just deduced \$6 \textbackslash{}neq \textbackslash{}text\{pos \} 4\$. So 6 is not Black.} \\
\texttt{    *   Therefore, **5 must be Black in Guess 3**. So **\$5 = \textbackslash{}text\{pos \} 2\$**.} \\
\texttt{*   From Guess 4 (`5 4 2 6`): 1B, 3W.} \\
\texttt{    *   Digits: 5, 4, 2, 6.} \\
\texttt{    *   Positions: 5@1, 4@2, 2@3, 6@4.} \\
\texttt{    *   We know \$2 = \textbackslash{}text\{pos \} 1\$ (so 2@3 is White).} \\
\texttt{    *   We know \$5 = \textbackslash{}text\{pos \} 2\$ (so 5@1 is White).} \\
\texttt{    *   We need 1 Black. It must be 4 or 6.} \\
\texttt{    *   If 4 is Black (\$4 = \textbackslash{}text\{pos \} 2\$): But we have \$5 = \textbackslash{}text\{pos \} 2\$. Contradiction.} \\
\texttt{    *   Therefore, **6 must be Black in Guess 4**. So **\$6 = \textbackslash{}text\{pos \} 4\$**.} \\
\texttt{    *   Wait, earlier I deduced \$6 \textbackslash{}neq \textbackslash{}text\{pos \} 4\$ from Guess 2?} \\
\texttt{        *   In Guess 2 (`2 1 5 6`), if 2 is Black, then 6 is White. So \$6 \textbackslash{}neq \textbackslash{}text\{pos \} 4\$.} \\
\texttt{        *   But in Guess 4, if 6 is Black, then \$6 = \textbackslash{}text\{pos \} 4\$.} \\
\texttt{        *   Contradiction.} \\
\\ 
\texttt{Let's re{-}evaluate "2 is Black in Guess 2".} \\
\texttt{Is it possible 5 is Black in Guess 2?} \\
\texttt{If \$5 = \textbackslash{}text\{pos \} 3\$:} \\
\texttt{*   Then in Guess 3 (`1 3 5 6`), 5 is at pos 2. So 5 is White.} \\
\texttt{*   Then 6 must be Black in Guess 3. So \$6 = \textbackslash{}text\{pos \} 4\$.} \\
\texttt{*   If \$6 = \textbackslash{}text\{pos \} 4\$:} \\
\texttt{    *   In Guess 2 (`2 1 5 6`), 6 is at pos 4 (Black).} \\
\texttt{    *   Then 2 and 5 are White.} \\
\texttt{    *   \$2 \textbackslash{}neq \textbackslash{}text\{pos \} 1\$.} \\
\texttt{    *   \$5 \textbackslash{}neq \textbackslash{}text\{pos \} 3\$. (But we assumed \$5 = \textbackslash{}text\{pos \} 3\$). Contradiction.} \\
\texttt{    *   So 5 cannot be Black in Guess 2.} \\
\\ 
\texttt{Is it possible 6 is Black in Guess 2?} \\
\texttt{If \$6 = \textbackslash{}text\{pos \} 4\$:} \\
\texttt{*   Then in Guess 3 (`1 3 5 6`), 6 is at pos 4 (Black).} \\
\texttt{*   Then 5 is White (\$5 \textbackslash{}neq \textbackslash{}text\{pos \} 2\$).} \\
\texttt{*   In Guess 4 (`5 4 2 6`), 6 is at pos 4 (Black).} \\
\texttt{*   Then 5, 4, 2 are White.} \\
\texttt{    *   \$5 \textbackslash{}neq \textbackslash{}text\{pos \} 1\$.} \\
\texttt{    *   \$4 \textbackslash{}neq \textbackslash{}text\{pos \} 2\$.} \\
\texttt{    *   \$2 \textbackslash{}neq \textbackslash{}text\{pos \} 3\$.} \\
\texttt{*   In Guess 5 (`6 5 2 4`), 6 is at pos 1.} \\
\texttt{    *   If 6 is Black in Guess 2 (\$6=\textbackslash{}text\{pos \} 4\$), then in Guess 5, 6 is at pos 1 (White).} \\
\texttt{    *   We need 1 Black in Guess 5.} \\
\texttt{    *   Candidates: 5, 2, 4.} \\
\texttt{    *   We know \$5 \textbackslash{}neq \textbackslash{}text\{pos \} 2\$ (from G3).} \\
\texttt{    *   We know \$4 \textbackslash{}neq \textbackslash{}text\{pos \} 4\$ (from G1).} \\
\texttt{    *   We know \$2 \textbackslash{}neq \textbackslash{}text\{pos \} 3\$ (from G4).} \\
\texttt{    *   Also \$2 \textbackslash{}neq \textbackslash{}text\{pos \} 2\$ (from G1).} \\
\texttt{    *   So 2 must be at pos 1.} \\
\texttt{    *   If \$2 = \textbackslash{}text\{pos \} 1\$:} \\
\texttt{        *   In Guess 5 (`6 5 2 4`), 2 is at pos 3. So 2 is White.} \\
\texttt{        *   So Black must be 5 or 4.} \\
\texttt{        *   If 5 is Black (\$5 = \textbackslash{}text\{pos \} 2\$): Contradicts \$5 \textbackslash{}neq \textbackslash{}text\{pos \} 2\$.} \\
\texttt{        *   If 4 is Black (\$4 = \textbackslash{}text\{pos \} 4\$): Contradicts \$4 \textbackslash{}neq \textbackslash{}text\{pos \} 4\$.} \\
\texttt{    *   So 6 cannot be Black in Guess 2.} \\
\\ 
\texttt{This implies my initial set of digits or the interpretation of "1B, 2W" is tricky.} \\
\texttt{Let's try a different combination.} \\
\texttt{Maybe the code is **4 2 6 5**?} \\
\texttt{G1 `1 2 3 4`: 2@2(W), 4@4(W). 0B, 2W. OK.} \\
\texttt{G2 `2 1 5 6`: 2@1(W), 5@3(W), 6@4(W). 0B, 3W. No.} \\
\\ 
\texttt{Maybe **2 6 5 4**?} \\
\texttt{G1 `1 2 3 4`: 2@2(W), 4@4(W). 0B, 2W. OK.} \\
\texttt{G2 `2 1 5 6`: 2@1(B), 5@3(W), 6@4(W). 1B, 2W. OK.} \\
\texttt{G3 `1 3 5 6`: 5@2(W), 6@4(W). 0B, 2W. No (Need 1B, 1W).} \\
\\ 
\texttt{Maybe **2 5 6 4**?} \\
\texttt{G1 `1 2 3 4`: 2@2(W), 4@4(W). 0B, 2W. OK.} \\
\texttt{G2 `2 1 5 6`: 2@1(B), 5@3(W), 6@4(W). 1B, 2W. OK.} \\
\texttt{G3 `1 3 5 6`: 5@2(B), 6@4(W). 1B, 1W. OK.} \\
\texttt{G4 `5 4 2 6`: 5@1(W), 4@2(W), 2@3(W), 6@4(W). 0B, 4W. No (Need 1B, 3W).} \\
\\ 
\texttt{Maybe **4 5 2 6**?} \\
\texttt{G1 `1 2 3 4`: 2@3(W), 4@4(W). 0B, 2W. OK.} \\
\texttt{G2 `2 1 5 6`: 5@3(B), 2@3(W? No, 2 is at 3 in code, 5 is at 3 in guess).} \\
\texttt{   Guess: 2, 1, 5, 6. Code: 4, 5, 2, 6.} \\
\texttt{   Pos 1: 2 vs 4.} \\
\texttt{   Pos 2: 1 vs 5.} \\
\texttt{   Pos 3: 5 vs 2.} \\
\texttt{   Pos 4: 6 vs 6 (B).} \\
\texttt{   Result: 1B, 3W. No (Need 1B, 2W).} \\
\\ 
\texttt{Let's try **5 2 6 4**.} \\
\texttt{G1 `1 2 3 4`: 2@2(W), 4@4(W). 0B, 2W. OK.} \\
\texttt{G2 `2 1 5 6`: 5@1(W), 6@4(W). 0B, 2W. No.} \\
\\ 
\texttt{Let's try **6 2 5 4**.} \\
\texttt{G1 `1 2 3 4`: 2@2(W), 4@4(W). 0B, 2W. OK.} \\
\texttt{G2 `2 1 5 6`: 5@3(W), 6@4(W). 0B, 2W. No.} \\
\\ 
\texttt{Let's try **2 4 6 5**.} \\
\texttt{G1 `1 2 3 4`: 2@2(W), 4@3(W). 0B, 2W. OK.} \\
\texttt{G2 `2 1 5 6`: 2@1(B), 5@4(W). 1B, 1W. No.} \\
\\ 
\texttt{Let's try **4 6 2 5**.} \\
\texttt{G1 `1 2 3 4`: 2@3(W), 4@1(W). 0B, 2W. OK.} \\
\texttt{G2 `2 1 5 6`: 5@4(W), 6@2(W). 0B, 2W. No.} \\
\\ 
\texttt{Let's try **5 6 2 4**.} \\
\texttt{G1 `1 2 3 4`: 2@3(W), 4@4(W). 0B, 2W. OK.} \\
\texttt{G2 `2 1 5 6`: 5@1(W), 6@2(W). 0B, 2W. No.} \\
\\ 
\texttt{Let's try **6 4 5 2**.} \\
\texttt{G1 `1 2 3 4`: 2@4(W), 4@2(W). 0B, 2W. OK.} \\
\texttt{G2 `2 1 5 6`: 5@3(B), 6@1(W). 1B, 1W. No.} \\
\\ 
\texttt{Let's try **4 2 5 6**.} \\
\texttt{G1 `1 2 3 4`: 2@2(W), 4@1(W). 0B, 2W. OK.} \\
\texttt{G2 `2 1 5 6`: 5@3(B), 6@4(B). 2B, 0W. No.} \\
\\ 
\texttt{Let's try **2 6 4 5**.} \\
\texttt{G1 `1 2 3 4`: 2@2(W), 4@3(W). 0B, 2W. OK.} \\
\texttt{G2 `2 1 5 6`: 2@1(B), 5@4(W). 1B, 1W. No.} \\
\\ 
\texttt{Let's try **5 4 6 2**.} \\
\texttt{G1 `1 2 3 4`: 2@4(W), 4@2(W). 0B, 2W. OK.} \\
\texttt{G2 `2 1 5 6`: 5@1(W), 6@3(W). 0B, 2W. No.} \\
\\ 
\texttt{Let's try **6 5 2 4** again.} \\
\texttt{G1 `1 2 3 4`: 2@3(W), 4@4(W). 0B, 2W. OK.} \\
\texttt{G2 `2 1 5 6`: 2@1(W), 5@3(W), 6@4(W). 0B, 3W. No.} \\
\\ 
\texttt{Wait, did I miss a permutation?} \\
\texttt{Let's try **4 5 6 2**.} \\
\texttt{G1 `1 2 3 4`: 2@4(W), 4@1(W). 0B, 2W. OK.} \\
\texttt{G2 `2 1 5 6`: 5@2(W), 6@3(W). 0B, 2W. No.} \\
\\ 
\texttt{Let's try **2 5 4 6**.} \\
\texttt{G1 `1 2 3 4`: 2@2(W), 4@3(W). 0B,} \\
            }
        }
    }
    & & \\ \\

    \theutterance \stepcounter{utterance}  
    & & & \multicolumn{2}{p{0.3\linewidth}}{
        \cellcolor[rgb]{0.9,0.9,0.9}{
            \makecell[{{p{\linewidth}}}]{
                \texttt{\tiny{[GM$|$GM]}}
                \texttt{role: Player 0} \\
\texttt{invalid\_move: True} \\
\texttt{turn\_count: 6} \\
\texttt{reason: Invalid Move: You have already guessed {[}2, 1, 5, 6{]}. Please try a different guess.} \\
\texttt{reward: 0.625} \\
            }
        }
    }
    & & \\ \\

\end{supertabular}
}

\end{document}
